% 
\section{Proof of \Cref{theorem: ord for the standard language valuation}}\label{section: theorem: ord for the standard language valuation}
\begin{definition}
       Let \[\SUBW_{\mathrm{st}} \defeq \bigcup_{n \ge 0} \left\{\val \in \SUBW_{n} \mathrel{\Big|}
              \begin{aligned}
                     \mbox{$\bigcup_{a \in \vsig} \val(a) = \set{\ell_{i} \mid i \in \range{0, n-1}}$} \\
                     \mbox{$\val(a)$ (where $a \in \vsig$) are disjoint sets}
              \end{aligned}
              \right\}.\]
\end{definition}

\begin{proposition}\label{proposition: st surjective}
Let $n \in \nat$ and let $\val' \in \SUBW_{n} \cap \SUBW_{\mathrm{st}}$.
For $i \in \range{0, n-1}$, let $a_i$ be the unique one s.t.\ $\val'(a_i) = \set{\ell_{i}}$.
Then $\val_{\mathrm{st}}^{\tuple{a_{0}, \dots, a_{n-1}}} = \val'$.
\end{proposition}
\begin{proof}
By \Cref{definition: word-to-letter valuations},
$\ell_{i} \in \val_{\mathrm{st}}^{\tuple{a_{0}, \dots, a_{n-1}}}(x)$ iff
$a_{i} \in \val_{\mathrm{st}}(x)$ iff $a_i = x$ iff
$\ell_{i} \in \val'(x)$.
Also, both $\val_{\mathrm{st}}^{\tuple{a_{0}, \dots, a_{n-1}}}(x)$ and $\val'(x)$ are subsets of $\set{\ell_{0}, \dots, \ell_{n-1}}$.
Thus, $\val_{\mathrm{st}}^{\tuple{a_{0}, \dots, a_{n-1}}} = \val'$.
\end{proof}

\begin{lemma}[cf.\ \Cref{lemma: LANG and SUBW}]\label{lemma: st lang and SUBW}
       For all \kl{quantifier-free formulas} $\fml$, we have:
       \[\val_{\mathrm{st}} \models \fml \quad\Leftrightarrow\quad \SUBW_{\mathrm{st}} \models \fml.\]
\end{lemma}
\begin{proof}
       As with \Cref{lemma: LANG and SUBW}, it suffices to prove when $\fml$ is of the form $(\bigvee_{k = 1}^{m} \term[1]_{k} \le \emp) \lor (\lnot \term[2] \le \emp)$.
      \begin{align*}
              & \SUBW_{\mathrm{st}} \not\models (\bigvee_{k = 1}^{m} \term[1]_{k} \le \emp) \lor (\lnot \term[2] \le \emp)                                                                                                                                    \\
              & \Rightarrow\quad \exists n, a_0, \dots, a_{n-1} \in \vsig,\ 
              \val_{\mathrm{st}}^{\tuple{a_{0}, \dots, a_{n-1}}} \not\models (\bigvee_{k = 1}^{m} \term[1]_{k} \le \emp) \lor (\lnot \term[2] \le \emp)                                                        \tag{\Cref{proposition: st surjective}}      \\
              & \Rightarrow\quad \val_{\mathrm{st}} \not\models (\bigvee_{k = 1}^{m} \term[1]_{k} \le \emp) \lor (\lnot \term[2] \le \emp)        \tag{By \Cref{proposition: letter-to-letter valuations}}                                         \\
              & \Rightarrow\quad \exists n \in \nat, i_1, j_1, \dots, i_m, j_m \mbox{ s.t.\ $0 \le i_1 \le j_1 \le i_2 \le j_2 \le \dots \le i_m \le j_m \le n$},                                                                                    \\
              & \phantom{\Rightarrow\quad} \exists a_{0}, \dots, a_{n-1} \in \vsig,\ (\bigwedge_{k = 1}^{m}  a_{i_k} \dots a_{j_k-1} \in \hat{\val}_{\mathrm{st}}(\term[1]_k)) \land \hat{\val}_{\mathrm{st}}(\term[2]) = \emptyset                   \\
              & \Rightarrow\quad \SUBW_{\mathrm{st}} \not\models (\bigvee_{k = 1}^{m} \term[1]_{k} \le \emp) \lor (\lnot \term[2] \le \emp). \tag*{(By \Cref{proposition: letter-to-letter valuations} and $\val_{\mathrm{st}}^{\tuple{a_0, \dots, a_{n-1}}} \in \SUBW_{n} \cap \SUBW_{\mathrm{st}}$) \qedhere}
       \end{align*}
\end{proof}

\begin{proposition}[Cor.\ of \Cref{proposition: subw and subwrel iso}]\label{proposition: st SUBW and ORD}
       For all \kl{quantifier-free formulas} $\fml$, we have:
       \[\SUBW_{\mathrm{st}} \models \fml \quad\Leftrightarrow\quad \RSUB_{\mathrm{st}} \models \fml.\]
\end{proposition}
\begin{proof}
By the same bijective map $f \colon L \mapsto \set{\tuple{i,j} \mid 0 \le i \le j \le n \land \ell_{i} \dots \ell_{j-1} \in L}$ as \Cref{proposition: subw and subwrel iso}.
% where $\ell_{i} \dots \ell_{j-1}$ denotes $\id$ if $i = j$.
\end{proof}

Combining \Cref{lemma: st lang and SUBW,proposition: st SUBW and ORD} yields
\[\val_{\mathrm{st}} \models \fml \quad\Leftrightarrow\quad \RSUB_{\mathrm{st}} \models \fml.\]
Particularly, we have $\ljump{\term[1]} = \ljump{\term[2]} \Leftrightarrow \RSUB_{\mathrm{st}} \models \term[1] = \term[2]$ (\Cref{theorem: ord for the standard language valuation}).


\section{Slight Extensions of \Cref{theorem: LANG and ORD}}\label{section: LANG and ORD for positive quantifier-free formulas}
In this section, we note that we can extend \Cref{theorem: LANG and ORD} in the following two:
\begin{theorem}\label{theorem: LANG and ORD for positive quantifier-free formulas}
    For all \kl[positive quantifier-free formulas]{\emph{positive} quantifier-free formulas} $\fml$ of $\KA_{\set{-}}$ \kl{terms},
    we have: $\LANG \models \fml \Rightarrow \RSUB \models \fml$.
\end{theorem}
\begin{proof}[Proof Sketch]
    By the same surjective $\sig$-homomorphism in the proof of \Cref{theorem: LANG and ORD}($\Rightarrow$).
\end{proof}
\begin{theorem}\label{theorem: LANG and ORD for quantifier-free formulas}
    For all \kl{quantifier-free formulas} $\fml$ of $\KA_{\set{-}}$ \kl{terms},
    we have: $\RSUB \models \fml \Rightarrow \LANG \models \fml$.
\end{theorem}
\begin{proof}
    Because the \kl{formulas} $\term[1] = \term[2] \leftrightarrow (\term[1] \le \term[2] \land \term[2] \le \term[1])$ and $\term[1] \le \term[2] \leftrightarrow \term[1] \cap \term[2]^{-} \le \emp$ are valid on $\LANG \cup \SUBW$,
    without loss of generality, we can assume that each \kl{equation} in $\fml$ is of the form $\term[3] \le \emp$.
    By taking the conjunctive normal form,
    it suffices to prove when $\fml$ is of the form $(\bigvee_{i = 1}^{n} \term[1]_i \le \emp) \lor (\bigvee_{j = 1}^{m} \lnot \term[2]_j \le \emp)$.
    We prove the contraposition.
    By $\LANG \not\models \fml$, there are $X, \val \in \LANG_X, \word_1, \dots, \word_n \in X^*$ such that
    $\word_i \in \hat{\val}(\term[1]_i)$ for $i \in \range{1,n}$ and $\hat{\val}(\term[2]_j) = \emptyset$ for $j \in \range{1,m}$.
    By letting $\word_0 \defeq \word_1 \dots \word_n$ and considering the same $\sig$-homomorphism as \Cref{theorem: LANG and ORD}($\Leftarrow$), we have $\RSUB \not\models \fml$.
\end{proof}
For \Cref{theorem: LANG and ORD for quantifier-free formulas},
note that the converse direction fails (\Cref{remark: quantifier-free}), cf.\ \Cref{theorem: LANG and ORD for positive quantifier-free formulas}.

\section{Proof of \Cref{corollary: graph characterization}}\label{section: corollary: graph characterization}
\begin{proof}
    We have:
    \begin{align*}
        &\algclass \models \term[1] \le \term[2] \ \Leftrightarrow\  \forall \word \in \ljump{\term[1]}_{\tilde{\vsig}}, \forall \graph[2] \in \FillS_{\algclass}(\const{\graph}(\word)), \exists \word[2] \in \ljump{\term[2]}_{\tilde{\vsig}}, \const{\graph}(\word[2]) \homo \graph[2]^{\FillQ} \tag{\Cref{theorem: graph characterization} and $\glang(\term[2]) = \set{\const{\graph}(\word[2]) \mid \word[2] \in \ljump{\term[2]}_{\tilde{\vsig}}}$} \\
                                                & \Leftrightarrow\quad \forall \word \in \ljump{\term[1]}_{\tilde{\vsig}}, \forall \graph[2] \in \FillS_{\algclass}(\const{\graph}(\word)), \ljump{\term[2]}_{\tilde{\vsig}} \cap \ljump{\graph[2]}_{\tilde{\vsig}} \neq \emptyset \tag{$\ljump{\graph[2]}_{\tilde{\vsig}} = \set{\word[2] \in \tilde{\vsig}^* \mid \const{\graph}(\word[2]) \homo \graph[2]^{\FillQ}}$}                       \\
                                                & \Leftrightarrow\quad \ljump{\term[1]}_{\tilde{\vsig}} \subseteq \set{\word \in \tilde{\vsig}^* \mid \forall \graph[2] \in \FillS_{\algclass}(\const{\graph}(\word)), \ljump{\term[2]}_{\tilde{\vsig}} \cap \ljump{\graph[2]}_{\tilde{\vsig}} \neq \emptyset }. \tag*{\qedhere}
    \end{align*}
\end{proof}

\section{Proof of \Cref{proposition: witness}}\label{section: proposition: witness}
Let $R' \subseteq \domain{\graph[2]} \times \domain{\graph[3]}$ be the minimal set such that
\begin{itemize}
       \item $\tuple{\src^{\graph[2]}, \src^{\graph[3]}} \in R'$,
       \item $\forall a \in \tilde{\vsig}_{\id}, \forall x, x' \in \domain{\graph[2]}, \forall y, y' \in \domain{\graph[3]}$,\   $(\tuple{x, y} \in R' \land \tuple{x, x'} \in \delta_{a}^{\graph[2]} \land \tuple{y, y'} \in \delta_{a}^{\graph[3]}) \Rightarrow \tuple{x', y'} \in R'$.
\end{itemize}
\begin{claim}\label{claim: witness pre}
$\ljump{\graph[3]}_{\tilde{\vsig}} \cap \ljump{\graph[2]}_{\tilde{\vsig}} \neq \emptyset$ \quad$\Leftrightarrow$\quad  $\tuple{\tgt^{\graph[2]}, \tgt^{\graph[3]}} \in R'$.
\end{claim}
\begin{claimproof}
By definition, $R'$ coincides with the set of all reachable states of the product \kl{NFA} of $\graph[2]$ and $\graph[3]$.
\end{claimproof}

Let $R \subseteq \domain{\graph[2]} \times \domain{\graph[3]}$ be the minimal set such that
\begin{itemize}
       \item $\tuple{\src^{\graph[2]}, \src^{\graph[3]}} \in R$,
       \item $\forall a \in \tilde{\vsig}_{\id}, \forall x, x' \in \domain{\graph[2]}, \forall y, y' \in \domain{\graph[3]}$,\   $(\tuple{x, y} \in R \land \tuple{x, x'} \in a^{\graph[2]} \land \tuple{y, y'} \in \delta_{a}^{\graph[3]}) \Rightarrow \tuple{x', y'} \in R$.
\end{itemize}
\begin{claim}\label{claim: witness eq}
$R = R'$.
\end{claim}
\begin{claimproof}
($\subseteq$):
Clear, by $a^{\graph[2]} \subseteq \delta_{a}^{\graph[2]}$.
($\supseteq$):
By induction on derivations of $R'$.
\begin{itemize}
\item Case $\tuple{\src^{\graph[2]}, \src^{\graph[3]}} \in R'$:
Trivial, by $\tuple{\src^{\graph[2]}, \src^{\graph[3]}} \in R$.
\item Case $(\tuple{x, y} \in R' \land \tuple{x, x'} \in \delta_{a}^{\graph[2]} \land \tuple{y, y'} \in \delta_{a}^{\graph[3]}) \Rightarrow \tuple{x', y'} \in R'$:
By IH, $\tuple{x, y} \in R$.
\begin{itemize}
       \item Sub-Case $a \neq \id$:
Let $x_{0}, \dots, x_{n-1}, x_{n}, \dots, x_{m}$ be s.t.\ $\tuple{x, x'} = \tuple{x_0, x_m}$ and 
\begin{itemize}
       \item for all $i \in \range{1, n-1}$,  $\tuple{x_{i-1}, x_{i}} \in \id^{\graph[2]}$,
       \item $\tuple{x_{n-1}, x_{n}} \in a^{\graph[2]}$,
       \item for all $i \in \range{n+1, m}$,  $\tuple{x_{i-1}, x_{i}} \in \id^{\graph[2]}$.
\end{itemize}
Let $y_0 = \dots = y_{n-1} = y$ and $y_{n} = \dots = y_{m} = y'$.
Then by applying the second rule multiply, we have $\tuple{x', y'} \in R$. 
       \item Sub-Case $a = \id$:
       By reflexivity of $\id^{\graph[2]}$, $\tuple{x, x'} \in (\id^{\graph[2]})^{+}$.
       Let $x_{0}, \dots, x_{m}$ ($m > 0$) be s.t.\ $\tuple{x, x'} = \tuple{x_0, x_m}$ and 
\begin{itemize}
       \item for all $i \in \range{1, m}$,  $\tuple{x_{i-1}, x_{i}} \in \id^{\graph[2]}$.
\end{itemize}
       Let $y_0 = y$ and $y_{1} = \dots = y_{m} = y'$.
       Then by applying the second rule multiply, we have $\tuple{x', y'} \in R$. 
       \claimqedhere
\end{itemize}
\end{itemize}
\end{claimproof}
       
\begin{proof}[Proof of \Cref{proposition: witness}]
       ($\Rightarrow$):
       By letting $U$ as the map defined by $U(x) \defeq \set{y \mid \tuple{x, y} \in R}$.
       Here, $\tgt^{\graph[3]} \not\in U_{\tgt^{\graph[2]}}$ is shown by $\ljump{\graph[3]}_{\tilde{\vsig}} \cap \ljump{\graph[2]}_{\tilde{\vsig}} = \emptyset$ with \Cref{claim: witness pre,claim: witness eq}.
       ($\Leftarrow$):
       Let $R'' \defeq \set{\tuple{x, y} \mid y \in U(x)}$.
       By the minimality of $R$, we have $R \subseteq R''$.
       By $\tuple{\tgt^{\graph[2]}, \tgt^{\graph[3]}} \not\in R''$, we have $\tuple{\tgt^{\graph[2]}, \tgt^{\graph[3]}} \not\in R$.
       Hence by \Cref{claim: witness pre,claim: witness eq}, we have $\ljump{\graph[3]}_{\tilde{\vsig}} \cap \ljump{\graph[2]}_{\tilde{\vsig}} = \emptyset$.
\end{proof}

% \begin{proof}
%        \ref{proposition: witness 1}$\Rightarrow$\ref{proposition: witness 2}:
%        Let
%        $U_{i} \defeq \set{z \in \domain{\automaton} \mid \exists \word[2] \in \ljump{\automaton[1][\bl, z]},\  \graph(\word[2]) \homo \graph[2]^{\FillQ}[\bl, \quo{i}_{\graph[2]}]}$.
%        For \ref{definition: witness s-t}:
%        Clear, by the assumption.
%        For \ref{definition: witness Con}:
%        Let $\tuple{i, j} \in a^{\graph[2]}, x \in U_{i}, y \in \delta_{a}^{\automaton}(x)$.
%        Then we show $y \in U_{j}$.
%        By $x \in U_{i}$, there is $\word[2] \in \ljump{\automaton[1][\bl, x]}$ s.t.\ $\graph(\word[2]) \homo \graph[2]^{\FillQ}[\bl,\quo{i}_{\graph[2]}]$.
%        Case $a \neq \id$:
%        Since $\word[2]a \in \ljump{\automaton[1][\bl, y]}$ and $\tuple{i, j} \in a^{\graph[2]}$,
%        we have that $\graph(\word[2] a) \homo \graph[2]^{\FillQ}[\bl, \quo{j}_{\graph[2]}]$; thus, $y \in U_{j}$.
%        Case $a = \id$:
%        Since $\quo{i}_{\graph[2]} = \quo{j}_{\graph[2]}$,
%        we have that $\graph(\word[2]) \homo \graph[2]^{\FillQ}[\bl, \quo{j}_{\graph[2]}]$; thus, $y \in U_{j}$.
%        \ref{proposition: witness 2}$\Rightarrow$\ref{proposition: witness 1}:
%        Assume, towards contradiction, that there is $\word \in \ljump{\automaton}$ s.t.\  $\graph(\word) \homo \graph[2]^{\FillQ}$.
%        Let $\word = a_1 \dots a_n$.
%        Let $q_0, \dots, q_n \in \domain{\automaton}$ be such that $\tuple{q_{i-1}, q_{i}} \in \delta_{a_i}^{\automaton}$ for all $i$ and $\tuple{q_0, q_n} = \tuple{\src^{\automaton}, \tgt^{\automaton}}$.
%        Let $v_0, v_0', v_1, \dots, v_n, v_n' \in \domain{\graph[2]}$ be such that $\tuple{v_{i-1}', v_{i}} \in a_i^{\graph[2]}$ and $\tuple{v_{i}', v_{i}} \in \id^{\graph[2]}$ for all $i$ and $\tuple{v_0, v_n'} = \tuple{\src^{\graph[2]}, \tgt^{\graph[2]}}$.
%        We prove that $q_i \in U_{v_i'}$ for all $i$.
%        Case $i = 0$:
%        By \ref{definition: witness s-t}, $q_0 \in U_{v_0} = U_{v_0'}$ (\Cref{proposition: con I}).
%        Case $i > 0$:
%        \begin{align*}
%               q_{i} \in \delta_{a}^{\automaton}(q_{i-1}) & \subseteq \delta_{a}^{\automaton}(U_{v_{i-1}'}) \tag{IH}                                           \\
%                                                          & \subseteq U_{v_i} = U_{v_i'}. \tag{By \ref{definition: witness Con} and \Cref{proposition: con I}}
%        \end{align*}
%        Hence we have $q_n \in U_{v_n'}$, which contradicts $\tgt^{\automaton} \not\in U_{\tgt^{\graph[2]}}$.
% \end{proof}
\section{Proof of \Cref{lemma: saturable path}}\label{section: lemma: saturable path}
\begin{proof}
($\Rightarrow$):
By \Cref{proposition: witness}, let $\graph[2]' \in \FillS_{\RSUB}(\const{\graph}(\word))$ and let $U$ be an \kl{emptiness-witness} for $\ljump{\graph[3]}_{\tilde{\vsig}} \cap \ljump{\graph[2]'}_{\tilde{\vsig}} = \emptyset$.
We define the \kl{graph} $\graph[2]$ as follows:
\begin{itemize}
       \item $\domain{\graph[2]} = \domain{\graph[2]'}$,
       \item $a^{\graph[2]} = a^{\graph[2]'}$ for $a \in \set{\top, \id, \com{\id}}$,
       \item $a^{\graph[2]} = a^{\const{\graph}(\word)} \cup (a^{\graph[2]'} \cap \id^{\graph[2]'})$ for $a \in \tilde{\vsig}_{\id} \setminus \set{\top, \id, \com{\id}}$.
\end{itemize}
% Let $\graph_0 \in \allgraph_{\RSUB}$ and $h \colon \const{\graph}(\word) \homo \graph_0$ be s.t.\ $\graph[2]' = \FillS(h)$.
We then have that the pair $P \defeq \tuple{\graph[2], U}$ is a \kl{saturable path} for $\word \in \const{L}_{\graph[3]}$, as follows:
\begin{itemize}
       \item \ref{definition: saturable path Ext}:
              By that $\graph[2]'$ is an \kl{edge-saturation} w.r.t.\ $\RSUB$.
       \item \ref{definition: saturable path Con}:
       Because $\graph[2]'$ is \kl{consistent} by $\graph[2]' \in \FillS_{\RSUB}(\const{\graph}(\word))$.
\item \ref{definition: saturable path Wit}:
              Because $U$ is an \kl{emptiness-witness} for $\ljump{\graph[3]}_{\tilde{\vsig}} \cap \ljump{\graph[2]'}_{\tilde{\vsig}} = \emptyset$.
       \item \ref{definition: saturable path Sat}:
              Because $a^{\graph[2]'} \cup \com{a}^{\graph[2]'} = \top^{\graph[2]'}$ and $U$ is an \kl{emptiness-witness} for $\ljump{\graph[3]}_{\tilde{\vsig}} \cap \ljump{\graph[2]'}_{\tilde{\vsig}} = \emptyset$.
\end{itemize}

% \ref{definition: saturable path Ext}:
% By the construction of $\graph[2]$.
% \ref{definition: saturable path Con}:
% Because $\graph[2]' \in \FillS_{\RSUB}(\const{\graph}(\word))$ is an \kl{edge extension} of $\graph[2]$, $(\graph[2]')$ is consistent.
% Thus $\graph[2]$ is consistent.
% \ref{definition: saturable path Wit}:
% Because $\graph[2]' \in \FillS_{\RSUB}(\const{\graph}(\word))$ is an \kl{edge extension} of $\graph[2]$ and $U$ is an \kl{emptiness-witness} for $\graph[3] \centernot\homo (\graph[2]')$,
% we have that $U$ is an \kl{emptiness-witness} for $\graph[3] \centernot\homo \graph[2]$ (\Cref{proposition: witness monotone}).
% \ref{definition: saturable path Sat}:
% Because $U$ is an \kl{emptiness-witness} for $\graph[3]_{\tilde{\vsig}} \centernot\homo (\graph[2]')$ and $x^{\graph[2]'} \cup \com{x}^{\graph[2]'} = \top^{\graph[2]'}$ for each $x \in \vsig$.
($\Leftarrow$):
Let $P = \tuple{\graph[2], U}$ be a \kl{saturable path} for $\word \in \const{L}_{\graph[3]}$.
By \ref{definition: saturable path Ext}, $\id^{\graph[2]}$ is an equivalence relation.
We define the \kl{graph} $\graph[2]'$ as follows:
\begin{itemize}
       \item $\univ{\graph[2]'} = \univ{\graph[2]}$,
       \item $a^{\graph[2]'} = a^{\graph[2]}$ for $a \in \set{\top, \id, \com{\id}}$,
       \item for $a \in \vsig$ and $\tuple{x, y} \in \top^{\graph[2]}$,
       \begin{itemize}
              \item if $\tuple{\quo{x}_{\id^{\graph[2]}}, \quo{y}_{\id^{\graph[2]}}} \in a^{\graph[2]^{\FillQ}}$, then $\tuple{x, y} \in a^{\graph[2]'} \setminus \com{a}^{\graph[2]'}$,
              \item else if $\tuple{\quo{x}_{\id^{\graph[2]}}, \quo{y}_{\id^{\graph[2]}}} \in \com{a}^{\graph[2]^{\FillQ}}$, then $\tuple{x, y} \in \com{a}^{\graph[2]'} \setminus a^{\graph[2]'}$,
              \item else if $U_{y} \subseteq \delta_{a}^{\graph[3]}(U_{x})$, then $\tuple{x, y} \in a^{\graph[2]'} \setminus \com{a}^{\graph[2]'}$,
              \item else $\tuple{x, y} \in \com{a}^{\graph[2]'} \setminus a^{\graph[2]'}$.
       \end{itemize}
\end{itemize}
By the construction of $\graph[2]'$, we have the following:
\begin{itemize}
       \item $\graph[2]'$ is an \kl{edge-extension} of $\graph[2]$:
       By \ref{definition: saturable path Con}, if $\tuple{\quo{x}_{\id^{\graph[2]}}, \quo{y}_{\id^{\graph[2]}}} \in \com{a}^{\graph[2]^{\FillQ}}$, then $\tuple{\quo{x}_{\id^{\graph[2]}}, \quo{y}_{\id^{\graph[2]}}} \not\in a^{\graph[2]^{\FillQ}}$.
       \item $\graph[2]'$ is \kl{consistent}:
       If $\quo{x}_{\id^{\graph[2]}} = \quo{y}_{\id^{\graph[2]}}$ then $U_{x} = U_{y}$, because $U_{x} \subseteq \delta_{\id}^{\graph[3]}(U_{x}) \subseteq U_{y} \subseteq \delta_{\id}^{\graph[3]}(U_{y}) \subseteq U_{x}$ by \ref{definition: saturable path Wit}; thus, if $\quo{x}_{\id^{\graph[2]}} = \quo{x'}_{\id^{\graph[2]}}$ and $\quo{y}_{\id^{\graph[2]}} = \quo{y'}_{\id^{\graph[2]}}$, then $\tuple{x, y} \in a^{\graph[2]'}$ iff $\tuple{x', y'} \in a^{\graph[2]'}$.

       \item for $a \in \vsig$, $\com{a}^{\graph[2]'} = \top^{\graph[2]'} \setminus a^{\graph[2]'}$:
       Because $a^{\graph[2]'} \cup \com{a}^{\graph[2]'} = \top^{\graph[2]'}$ and $\graph[2]'$ is \kl{consistent}.
       
%     \item for $a \in \vsig$, $a^{\graph[2]'} \supseteq \id^{\graph[2]'}$ or $\com{a}^{\graph[2]'} \supseteq \id^{\graph[2]'}$:
%     By \ref{definition: saturable path Ext}.
\end{itemize}
From them and \ref{definition: saturable path Ext}, we have $\graph[2]' \in \FillS_{\RSUB}(\const{\graph}(\word))$.
Also, $U$ is an \kl{emptiness-witness} for $\ljump{\graph[3]}_{\tilde{\vsig}} \cap \ljump{\graph[2]'}_{\tilde{\vsig}} = \emptyset$ as follows.
For edges already in $\graph[2]$, it is shown by \ref{definition: saturable path Wit}.
For extended edges from $\graph[2]$, it is shown by the construction of $\graph[2]'$ (for the last case of the four cases above, by $U_{y} \not\subseteq \delta_{a}^{\graph[3]}(U_{x})$ and \ref{definition: saturable path Sat}, we have $U_{y} \subseteq \delta_{\com{a}}^{\graph[3]}(U_{x})$).
Hence, this completes the proof.
% Hence, this completes the proof.
% Then, we have:
% \begin{itemize}
%     \item $x^{\graph[2]'} \cup \com{x}^{\graph[2]'} = \top^{\graph[2]'}$, because the third or fourth condition above always holds by \Cref{definition: saturable path Sat};
%     \item $x^{(\graph[2]')^{\FillQ}} \cap \com{x}^{(\graph[2]')^{\FillQ}} = \emptyset$, by the construction of $\graph[2]'$ (for each $\tuple{i, i'}, \tuple{j, j'} \in \id^{\graph[2]'}$,
%           we have
%           that $\tuple{[i]_{\id^{\graph[2]}}, [j]_{\id^{\graph[2]}}} \in x^{\graph[2]^{\FillQ}}$ iff $\tuple{[i']_{\id^{\graph[2]}}, [j']_{\id^{\graph[2]}}} \in x^{\graph[2]^{\FillQ}}$
%           and have that $U_{j} \in \delta_{x}^{\automaton}(U_{i})$ iff $U_{j'} \in \delta_{x}^{\automaton}(U_{i'})$ (because $U_{i} = U_{i'}$ and $U_{i'} = U_{j'}$).
%           Thus, $\tuple{i, j} \in x^{\graph[2]'}$ iff $\tuple{i', j'} \in x^{\graph[2]'}$);
%     \item $\forall x \in \vsig,\  x^{\graph[2]'} \supseteq \id^{\graph[2]'} \lor \com{x}^{\graph[2]'} \supseteq \id^{\graph[2]'}$, by \ref{definition: saturable path Ext}.
% \end{itemize}$(\graph[2]')^{\FillQ}$ is \kl{isomorphic} to $\mathrm{G}(\val, 0, n)$ for some $n \in \nat$ and $\val \in \RSUB_{n}$:
\end{proof}


\section{Proof of \Cref{lemma: sound and completeness}}\label{section: lemma: sound and completeness}

\begin{proof}
       ($\subseteq$):
       Let $\word = a_1 \dots a_{n} \in \const{L}_{\graph[3]}$.
       Let $P = \tuple{\graph[2], U}$ be a \kl{saturable path} for $\word \in \const{L}_{\graph[3]}$.
       Let $\Aset \defeq \set{a \in \tilde{\vsig}_{\id} \mid a^{\graph[2]} \supseteq \id^{\graph[2]}}$ (note that $X \in \mathcal{X}$).
       For each $i$, let $\mathcal{U}_i \defeq \bigcup_{j; \tuple{j, i} \in \com{\id}^{\graph[2]}} U_{j}^2$.
       Then we have:
       \begin{itemize}
              \item $\fml^{\graph[3]}(\mathcal{U}_i, U_i)$: By \ref{definition: saturable path Sat} and \Cref{proposition: saturable}.
              \item $\forall a \in \Aset, \delta_{a}^{\graph[3]}(U_i) \subseteq U_i$:
                    By $a^{\graph[2]} \supseteq \id^{\graph[2]} \supseteq \Delta_{\domain{\graph[2]}}$ and \ref{definition: saturable path Wit}.
       \end{itemize}
       Thus $\tuple{\mathcal{U}_i, U_i} \in \domain{\graph[3]^{\FillS_{\Aset}}}$.
       We consider the following run of the \kl{NFA} $\graph[3]^{\FillS_{\Aset}}$ on $\word$:
       \[\begin{tikzpicture}[remember picture, baseline = -.7ex]
              \graph[grow right = 2.cm, branch down = 3.5ex, nodes={}]{
              {lx/{\scriptsize $\blacktriangleright$}}-!-
              {0x/{\scriptsize $\tuple{\mathcal{U}_0, U_0}$}}-!-
              {1x/{\scriptsize $\tuple{\mathcal{U}_1, U_1}$}}-!-
              {2x/{\scriptsize $\tuple{\mathcal{U}_2, U_2}$}}-!-
              {3x/{}[inner sep = 0, outer sep = 0]}-!-
              {4x/{\scriptsize $\tuple{\mathcal{U}_{n}, U_{n}}$}}-!-
              {rx/{\scriptsize $\blacktriangleleft$}}
              };
              \node[left = .5em of lx](l){};
              \node[right = .5em of rx](r){};
              \graph[use existing nodes, edges={color=black, pos = .5, earrow}, edge quotes={fill=white, inner sep=1pt,font= \scriptsize}]{
                  {lx} ->["$\id$"] {0x} ->["$a_1$"] {1x} ->["$a_2$"] {2x} --["\;$\dots$\;", pos = .8] {3x.center} ->["$a_{n}$"] {4x} ->["$\id$"] {rx};
                  l -> lx; rx -> r;
              }; 
          \end{tikzpicture}.\]
       This is indeed a run of the \kl{NFA} $\graph[3]^{\FillS_{\Aset}}$ as follows:
       \begin{itemize}
              \item $\tuple{\blacktriangleright, \tuple{\mathcal{U}_0, U_0}} \in \id^{\graph[3]^{\FillS_{\Aset}}}$:
                    By $\src^{\graph[3]} \in U_0$ \ref{definition: saturable path Wit} and $\mathcal{U}_0 = \emptyset$.
              \item $\tuple{\tuple{\mathcal{U}_n, U_n}, \blacktriangleleft} \in \id^{\graph[3]^{\FillS_{\Aset}}}$:
                    By $\tgt^{\graph[3]} \not\in U_{n}$ \ref{definition: saturable path Wit}.
              \item $\forall i \in \range{1, n}, \tuple{\tuple{\mathcal{U}_{i-1}, U_{i-1}}, \tuple{\mathcal{U}_{i}, U_{i}}} \in a_{i}^{\graph[3]^{\FillS_{\Aset}}}$:
                    We distinguish the following cases:
                    \begin{itemize}
                           \item Case $\tuple{i-1, i} \in \id^{\graph[2]}$:
                                   \begin{itemize}
                                          \item  $\mathcal{U}_{i} = \mathcal{U}_{i-1}$:
                                          By $\tuple{j, i} \in \com{\id}^{\graph[2]}$ iff $\tuple{j, i-1} \in \com{\id}^{\graph[2]}$, for all $j$.
                                          \item $U_{i} = U_{i-1}$:
                                          By \ref{definition: saturable path Wit}, we have $U_{i-1} \subseteq \delta_{\id}^{\graph[3]}(U_{i-1}) \subseteq U_{i} \subseteq \delta_{\id}^{\graph[3]}(U_{i}) \subseteq U_{i-1}$.
                                          \item $a_{i} \in \Aset$ ($a_{i}^{\graph[2]} \supseteq \id^{\graph[2]}$):
                                          By $a_{i}^{\graph[2]} \cap \id^{\graph[2]} \neq \emptyset$ and \ref{definition: saturable path Ext},
                                          we have $a_{i}^{\graph[2]} = a_{i}^{\const{\graph}(\word)} \cup \id^{\graph[2]}$ (if not, this contradicts to \ref{definition: saturable path Con}).
                                   \end{itemize}
                                 Thus by $\fml[2]_{a_{i}, \id}^{\Aset}(\mathcal{U}_{i-1}, U_{i-1}, \mathcal{U}_{i}, U_{i})$, we have $\tuple{\tuple{\mathcal{U}_{i-1}, U_{i-1}}, \tuple{\mathcal{U}_{i}, U_{i}}} \in a_{i}^{\graph[3]^{\FillS_{\Aset}}}$.
                           \item Case $\tuple{i-1, i} \in \com{\id}^{\graph[2]}$:
                           \begin{itemize}
                                   \item $\mathcal{U}_{i} = \mathcal{U}_{i-1} \cup U_{i-1}^2$:
                                   By $\tuple{j, i} \in \com{\id}^{\graph[2]}$ iff $j < i$ iff $\tuple{j, i-1} \in \com{\id}^{\graph[2]} \lor \tuple{j, i-1} \in \id^{\graph[2]}$, for all $j$.
                                   (Intuitively, $\mathcal{U}_{i-1}$ corresponds to the case $\tuple{j, i-1} \in \com{\id}^{\graph[2]}$ and $U_{i-1}^2$ corresponds to the case $\tuple{j, i-1} \in \id^{\graph[2]}$.)
                                   \item $\delta_{a_{i}}^{\graph[3]}(U_{i-1}) \subseteq U_{i}$:
                                   By \ref{definition: saturable path Wit}.
                                   \item $\delta_{\top}^{\graph[3]}(\set{u \mid \tuple{u, u} \in \mathcal{U}_{i}}) \subseteq U_{i}$:
                                   We have
                                   $\delta_{\top}^{\graph[3]}(\set{u \mid \tuple{u, u} \in \mathcal{U}_{i}})
                                   = \delta_{\top}^{\graph[3]}(\bigcup_{j; \tuple{j, i} \in \com{\id}^{\graph[2]}} U_{j})
                                   = \bigcup_{j < i} \delta_{\top}^{\graph[3]}(U_{j}) \subseteq U_{i}$ by \ref{definition: saturable path Wit}.
                                   \item $\delta_{\com{\id}}^{\graph[3]}(\set{u \mid \tuple{u, u} \in \mathcal{U}_{i}}) \subseteq U_{i}$:
                                   We have
                                   $\delta_{\com{\id}}^{\graph[3]}(\set{u \mid \tuple{u, u} \in \mathcal{U}_{i}})
                                   = \delta_{\com{\id}}^{\graph[3]}(\bigcup_{j; \tuple{j, i} \in \com{\id}^{\graph[2]}} U_{j})
                                   = \bigcup_{j < i} \delta_{\com{\id}}^{\graph[3]}(U_{j}) \subseteq U_{i}$ by \ref{definition: saturable path Wit}.
                           \end{itemize}
                                 Thus by $\fml[2]_{a_{i}, \com{\id}}^{\Aset}(\mathcal{U}_{i-1}, U_{i-1}, \mathcal{U}_{i}, U_{i})$, we have $\tuple{\tuple{\mathcal{U}_{i-1}, U_{i-1}}, \tuple{\mathcal{U}_{i}, U_{i}}} \in a_{i}^{\graph[3]^{\FillS_{\Aset}}}$.
                    \end{itemize}
       \end{itemize}
       Hence, $\word \in \ljump{\graph[3]^{\FillS_{\Aset}}}$.

       ($\supseteq$):
       Let $X \subseteq \mathcal{X}$ and $\word = a_1 \dots a_n \in \ljump{\graph[3]^{\FillS_{X}}}_{\tilde{\vsig}}$.
       Let the run of $\graph[3]^{\FillS_{X}}$ on $\word$ be as follows:
       \[\begin{tikzpicture}[remember picture, baseline = -.7ex]
              \graph[grow right = 2.cm, branch down = 3.5ex, nodes={}]{
              {lx/{\scriptsize $\blacktriangleright$}}-!-
              {0x/{\scriptsize $\tuple{\mathcal{U}_0, U_0}$}}-!-
              {1x/{\scriptsize $\tuple{\mathcal{U}_1, U_1}$}}-!-
              {2x/{\scriptsize $\tuple{\mathcal{U}_2, U_2}$}}-!-
              {3x/{}[inner sep = 0, outer sep = 0]}-!-
              {4x/{\scriptsize $\tuple{\mathcal{U}_{n}, U_{n}}$}}-!-
              {rx/{\scriptsize $\blacktriangleleft$}}
              };
              \node[left = .5em of lx](l){};
              \node[right = .5em of rx](r){};
              \graph[use existing nodes, edges={color=black, pos = .5, earrow}, edge quotes={fill=white, inner sep=1pt,font= \scriptsize}]{
                  {lx} ->["$\id$"] {0x} ->["$a_1$"] {1x} ->["$a_2$"] {2x} --["\;$\dots$\;", pos = .8] {3x.center} ->["$a_{n}$"] {4x} ->["$\id$"] {rx};
                  l -> lx; rx -> r;
              }; 
          \end{tikzpicture}.\]
       Let $\graph[2]$ be the \kl{edge-extension} of $\const{\graph}(\word)$ defined as follows:
       \begin{itemize}
              \item $\top^{\graph[2]} = \set{\tuple{x, y} \in \range{0, n}^2 \mid \forall i \in \range{y + 1, x},\  \lnot \fml[2]_{a_{i}, \com{\id}}^{\Aset}(\mathcal{U}_{i-1}, U_{i-1}, \mathcal{U}_{i}, U_{i})}$,
              \item $\id^{\graph[2]} = \top^{\graph[2]} \cap \set{\tuple{x, y} \mid \tuple{y, x} \in \top^{\graph[2]}}$ and $\com{\id}^{\graph[2]} = \top^{\graph[2]} \setminus \id^{\graph[2]}$,
              \item $\forall a \in \vsig \cap \Aset$, $\tuple{a^{\graph[2]}, \com{a}^{\graph[2]}} = \tuple{a^{\const{\graph}(\word)} \cup \id^{\graph[2]}, \com{a}^{\const{\graph}(\word)}}$.
       \end{itemize}
       Note that by definition of $\top^{\graph[2]}$, we have
       \begin{itemize}
         \item $\top^{\graph[2]} \supseteq \set{\tuple{x, y} \mid x \le y}$,
         \item $\top^{\graph[2]}$ is transitive (by case analysis).
       \end{itemize}
       Hence, $\top^{\graph[2]}$ is a total preorder and each equivalence class w.r.t.\ $\id^{\graph[2]}$ is an interval $\range{l, r}$.

       Let $P \defeq \tuple{\graph[2], U}$ where $U$ is defined as $i \mapsto U_i$ for $i \in \range{0, n}$.
       The following depicts $P$.
       \[\begin{tikzpicture}[remember picture, baseline = -.5ex]
              \graph[grow right = 1.5cm, branch down = 3.5ex, nodes={}]{
              {0x/{}[mynode]}-!-
              {1x/{}[mynode]}-!-
              {2x/{}[mynode]}-!-
              {3x/{}[inner sep = 0, outer sep = 0]}-!-
              {4x/{}[mynode, xshift = -4em]}-!-
              {5x/{}[mynode]}-!-
              {6x/{}[inner sep = 0, outer sep = 0]}-!-
              {7x/{}[mynode]}
              };
              \node[left = .5em of 0x](l){};
              \node[right = .5em of 7x](r){};
              \graph[use existing nodes, edges={color=black, pos = .5, earrow}, edge quotes={fill=white, inner sep=1pt,font= \scriptsize}]{
                  {0x, 1x, 2x, 4x, 5x, 7x} ->["$X$", out=60, in=120, looseness=18] {0x, 1x, 2x, 4x, 5x, 7x};
                  {0x} ->["$a_1$"] {1x};
              {1x} ->["$a_2$"] {2x};
              {2x} ->["\;$\dots$\;"] {4x};
              {4x} ->["$a_{i}$"] {5x};
              {4x} ->[bend left = 25] {5x};
              {5x} --["\;$\dots$\;"] {6x.center};
              {6x.center} ->["$a_{n}$"] {7x};
              l -> 0x; 7x -> r;
              }; 
              \path (4x) edge[opacity = 0, bend left = 25] node(4x-5x){} (5x);
              \node[above right = .5em and 0.em of 4x-5x, fill = white, inner sep = 0](4x-5x-p){\scriptsize $\begin{cases} \id & (\mbox{if $\lnot \fml[2]_{a_{i}, \com{\id}}^{\Aset}(\mathcal{U}_{i-1}, U_{i-1}, \mathcal{U}_{i}, U_{i})$})\\
                     \com{\id} & (\mbox{otherwise})
              \end{cases}$};
              \draw[color = gray] (4x-5x) -- (4x-5x-p.west);
              \node[below = .5ex of 0x]{\scriptsize $U_{0}$};
              \node[below = .5ex of 1x]{\scriptsize $U_{1}$};
              \node[below = .5ex of 2x]{\scriptsize $U_{2}$};
              \node[below = .5ex of 4x]{\scriptsize $U_{i-1}$};
              \node[below = .5ex of 5x]{\scriptsize $U_{i}$};
              \node[below = .5ex of 7x]{\scriptsize $U_{n}$};
          \end{tikzpicture}.\]
       Then $P$ is a \kl{saturable path} for $\word \in \const{L}_{\graph[3]}$ as follows:
       \begin{itemize}
              \item \ref{definition: saturable path Ext}:
                    By the definition of $\graph[2]$.

              \item \ref{definition: saturable path Con}:
                    Assume that $a^{\graph[2]^{\FillQ}} \cap \com{a}^{\graph[2]^{\FillQ}} \neq \emptyset$.
                    Let $x, x', y, y'$ be s.t.\
                    $\quo{x}_{\id^{\graph[2]}} = \quo{x'}_{\id^{\graph[2]}}$,
                    $\quo{y}_{\id^{\graph[2]}} = \quo{y'}_{\id^{\graph[2]}}$,
                    $\tuple{x, y} \in a^{\graph[2]}$, and $\tuple{x', y'} \in \com{a}^{\graph[2]}$.
                    WLOG, we can assume that $a \in \Aset$ and $\com{a} \not\in \Aset$.
                    Then, we have the following:
                    \begin{itemize}
                     \item $\tuple{x', y'} \in \com{a}^{\const{\graph}(\word)}$ (so, $x' = y' - 1$ and $a_{y'} = \com{a}$): By $\com{a}^{\graph[2]} = \com{a}^{\const{\graph}(\word)}$ (since $\com{a} \not\in \Aset$).

                     \item $\tuple{x, y} \in a^{\const{\graph}(\word)}$ (so, $x = y - 1$ and $a_{y} = a$):
                     If not, then by $a^{\graph[2]} = a^{\const{\graph}(\word)} \cup \id^{\graph[2]}$, we have $\quo{x}_{\id^{\graph[2]}} = \quo{y}_{\id^{\graph[2]}}$.
                     Thus, $\tuple{y', y' - 1} \in \id^{\graph[2]} (\subseteq \top^{\graph[2]})$.
                     By the definition of $\top^{\graph[2]}$, we have $\lnot \fml[2]_{a_{y'}, \com{\id}}^{\Aset}(\mathcal{U}_{y'-1}, U_{y'-1}, \mathcal{U}_{y'}, U_{y'})$.
                     By the definition of $a^{\graph[3]^{\FillS_{X}}}$, we have $\fml[2]_{a_{y'}, \id}^{\Aset}(\mathcal{U}_{y'-1}, U_{y'-1}, \mathcal{U}_{y'}, U_{y'})$,
                     so $\com{a} \in \Aset$.
                     This contradicts $\com{a} \not\in \Aset$.
                  
                     \item $(\range{x, x'} \cup \range{x', x}) \cap (\range{y, y'} \cup \range{y', y}) = \emptyset$ (so, $x = x'$ and $y = y'$):
                     If not, then because the interval between $x$ and $x'$ and that between $y$ and $y'$ have an intersection,
                     we have $\quo{x}_{\id^{\graph[2]}} = \quo{y}_{\id^{\graph[2]}}$.
                     Then, in the same manner as above, we have $\com{a} \in X$.
                     This contradicts $\com{a} \not\in \Aset$.
                    \end{itemize}
                     Thus, we reach a contradiction, because $a = a_{y} = a_{y'} = \com{a}$ (by $y = y'$).
                     Hence, $a^{\graph[2]^{\FillQ}} \cap \com{a}^{\graph[2]^{\FillQ}} = \emptyset$.

              \item \ref{definition: saturable path Sat}:
                     By the form of $\graph[3]^{\FillS_{\Aset}}$,
                     we have $\mathcal{U}_{x} = \begin{cases}
                                   \mathcal{U}_{x-1}                & (\tuple{x-1, x} \in \id^{\graph[2]})       \\
                                   \mathcal{U}_{x-1} \cup U_{x-1}^2 & (\tuple{x-1, x} \in \com{\id}^{\graph[2]})
                            \end{cases}$.
                     Thus, $\mathcal{U}_y = \bigcup_{x; \tuple{x, y} \in \com{\id}^{\graph[2]}} U_{x}^2$ ($\bigstar$).
                     By \Cref{proposition: saturable}, this completes the proof.
 
              \item \ref{definition: saturable path Wit}:
                    For $\src^{\graph[3]} \in U_0$ and $\tgt^{\graph[3]} \not\in U_n$, they are shown by the form of $\graph[3]^{\FillS_{\Aset}}$.
                    For $\forall a \in \tilde{\vsig}_{\id}, \forall \tuple{x, y} \in a^{\graph[2]}, \delta_{a}^{\graph[3]}(U_{x}) \subseteq U_{y}$,
                    we distinguish the following cases:
                    \begin{itemize}
                     \item Case $a = \id$:
                     Then we have 
                     \begin{itemize}
                            \item $U_{x} = U_{y}$:
                            By $\tuple{x, y} \in \id^{\graph[2]}$ and the form of $\graph[3]^{\FillS_{\Aset}}$, we have the following: $\forall z \in \range{y + 1, x},\ \fml[2]_{a_{z}, \id}^{\Aset}(\mathcal{U}_{z-1}, U_{z-1}, \mathcal{U}_{z}, U_{z})$.
                            Thus, $U_{y} = U_{y + 1} = \dots = U_{x}$.
                            \item $\delta_{\id}^{\graph[3]}(U_{x}) \subseteq U_{x}$:
                            By $\tuple{\mathcal{U}_{x}, U_{x}} \in \domain{\graph[3]^{\FillS_{\Aset}}}$.
                     \end{itemize}
                     Hence, $\delta_{\id}^{\graph[3]}(U_{x}) \subseteq U_{y}$.

               \item Case $a = \com{\id}$:
                     Let $z \in \range{x + 1, y}$ be such that
                     $\fml[2]_{a_{z}, \com{\id}}^{\Aset}(\mathcal{U}_{z - 1}, U_{z - 1}, \mathcal{U}_{z} , U_{z})$ and
                     $\forall z' \in \range{z+1, y}, \lnot \fml[2]_{a_{z'}, \com{\id}}^{\Aset}(\mathcal{U}_{z' - 1}, U_{z' - 1}, \mathcal{U}_{z'} , U_{z'})$.
                     Then we have
                     \begin{align*}
                            \delta_{\com{\id}}^{\graph[3]}(U_{x}) &\subseteq \delta_{\com{\id}}^{\graph[3]}(\set{u \mid \tuple{u, u} \in \mathcal{U}_{z}}) \tag{by ($\bigstar$) and $\tuple{x, z} \in \com{\id}^{\graph[2]}$ (by $\tuple{z -1, z} \in \com{\id}^{\graph[2]}$)}\\
                            &\subseteq U_{z} \tag{by $\fml[2]_{a_{z}, \com{\id}}^{\Aset}(\mathcal{U}_{z - 1}, U_{z - 1}, \mathcal{U}_{z} , U_{z})$}\\
                            &\subseteq U_{z+1} = \dots = U_{y}. \tag{by the form of $\graph[3]^{\FillS_{\Aset}}$, $\fml[2]_{a_{z'}, \id}^{\Aset}(\mathcal{U}_{z' - 1}, U_{z' - 1}, \mathcal{U}_{z'} , U_{z'})$}
                     \end{align*}

               \item Case $a = \top$: We distinguish the following two sub-cases:
                     \begin{itemize}
                     \item Case $\tuple{x, y} \in \com{\id}^{\graph[2]}$:
                     By the similar argument as Case $a = \com{\id}$.
                     \item Case $\tuple{x, y} \in \id^{\graph[2]}$:
                     By the similar argument as Case $a = \id$, we have $U_{x} = U_{y}$ and $\delta_{\top}^{\graph[3]}(U_{x}) \subseteq U_{x}$, and thus $\delta_{\top}^{\graph[3]}(U_{x}) \subseteq U_{y}$.
                     \end{itemize}

              \item Case $a \in \set{a, \com{a} \mid a \in \vsig}$: We distinguish the following sub-cases:
                     \begin{itemize}
                            \item Case $\tuple{x, y} \in \com{\id}^{\graph[2]}$:
                            By $\tuple{x, y} \in a^{\graph[2]} \cap \com{\id}^{\graph[2]} = a^{\const{\graph}(\word)}$, we have $x = y - 1$ and $a_{y} = a$.
                            Thus by $\fml[2]_{a_{y}, \com{\id}}^{\Aset}(\mathcal{U}_{y - 1}, U_{y - 1}, \mathcal{U}_{y} , U_{y})$, we have $\delta_{a}^{\graph[3]}(U_{x}) \subseteq U_{y}$.
                            \item Case $a \not\in \Aset$:
                            By $a^{\graph[2]} = a^{\const{\graph}(\word)}$, we have $x = y - 1$ and $a_{y} = a$.
                            By the form of $\graph[3]^{\FillS_{\Aset}}$ with $\lnot \fml[2]_{a_{y}, \id}^{\Aset}(\mathcal{U}_{y - 1}, U_{y - 1}, \mathcal{U}_{y} , U_{y})$ (since $a_{y} \not\in \Aset$),
                            we have $\fml[2]_{a_{y}, \com{\id}}^{\Aset}(\mathcal{U}_{y - 1}, U_{y - 1}, \mathcal{U}_{y} , U_{y})$.
                            Hence, $\delta_{a}^{\graph[3]}(U_{x}) \subseteq U_{y}$.
                            \item Case $\tuple{x, y} \in \id^{\graph[2]}$ and $a \in \Aset$:
                            By the similar argument as Case $a = \id$, we have $U_{x} = U_{y}$ (by $\tuple{x, y} \in \id^{\graph[2]}$) and $\delta_{a}^{\graph[3]}(U_{x}) \subseteq U_{x}$ (by $a \in \Aset$).
                            Thus, $\delta_{a}^{\graph[3]}(U_{x}) \subseteq U_{y}$. \qedhere
                     \end{itemize}
                     \end{itemize}
       \end{itemize}
\end{proof}

\ifiscameraready
\else

\newpage
\ 
\newpage

(Extra Appendices)
 % By $\val \models \term[1] = \term[2] \leftrightarrow (\term[1] \le \term[2] \land \term[2] \le \term[1])$ and $\val \models \term[1] \le \term[2] \leftrightarrow \term[1] \cap \term[2]^{-} \le \emp$ (which are set-theoretic equivalences) for $\val \in \LANG \cup \SUBW$, without loss of generality, we can assume that each \kl{equation} in $\fml$ is of the form $\term[3] \le \emp$.
    % By taking the conjunctive normal form, it suffices to show the claimed lemma when $\fml$ is of the form $(\bigvee_{k} \term[1]_{k} \le \emp) \lor (\bigvee_{k} \lnot \term[2]_{k} \le \emp)$.
    % By $\val \models (\term[1] \le \emp \land \term[2] \le \emp) \leftrightarrow \term[1] \cup \term[2] \le \emp$ (and $\val \models \mathsf{f} \leftrightarrow \top \le \emp$) for $\val \in \LANG \cup \SUBW$,
    % it suffices to show when $\fml$ is of the form $(\bigvee_{k = 1}^{m} \term[1]_{k} \le \emp) \lor (\lnot \term[2] \le \emp)$.
    % Then, we have:
    % $\SUBW \not\models (\bigvee_{k = 1}^{m} \term[1]_{k} \le \emp) \lor (\lnot \term[2] \le \emp)$
    % $\Rightarrow$ $\LANG \not\models (\bigvee_{k = 1}^{m} \term[1]_{k} \le \emp) \lor (\lnot \term[2] \le \emp)$ (by \Cref{proposition: surjective,proposition: letter-to-letter valuations})
    % $\Rightarrow$ $\exists n \in \nat, i_1, j_1, \dots, i_m, j_m \mbox{ s.t.\ $0 \le i_1 \le j_1 \le i_2 \le j_2 \le \dots \le i_m \le j_m \le n$}, \exists X, \val \in \LANG_{X}, a_{0}, \dots, a_{n-1} \in X,\ (\bigwedge_{k = 1}^{m}  a_{i_k} \dots a_{j_k-1} \in \hat{\val}(\term[1]_k)) \land \hat{\val}(\term[2]) = \emptyset$
    % $\Rightarrow$ $\SUBW \not\models (\bigvee_{k = 1}^{m} \term[1]_{k} \le \emp) \lor (\lnot \term[2] \le \emp)$ (by \Cref{proposition: letter-to-letter valuations} and $\val^{\tuple{a_0, \dots, a_{n-1}}} \in \SUBW_{n}$).
    % By \Cref{lemma: LANG and SUBW} and \Cref{proposition: subw and subwrel iso}.


\section{$\SUBW$: subword models}\label{section: SUBW}
In this section, we also define the class $\SUBW$ of \emph{\kl{subword models}}, which is isomorphic to $\RSUB$.
Apart from $\RSUB$, intuitively, this model is based on the following property in \kl{the standard language valuation} ($\val_{\mathrm{st}}$):
the membership $\word \in \hat{\val}_{\mathrm{st}}(\term)$ can be determined from $\word' \in \hat{\val}_{\mathrm{st}}(\term')$ where $\word'$ ranges over \emph{\kl{subwords}} of $\word$ (and $\term'$ ranges over subterms of $\term$).
This situation is the same also for any \kl{valuations} $\val \in \LANG$.
The class $\SUBW$ is helpful to use this property.
These models are \kl{language models} where the universe is restricted to the set of \kl{subwords} of a \kl{word}.

\begin{definition}\label{definition: subword language algebra}
    Let $X$ be a set and let $\word \in X^*$.
    The \intro*\kl{subword language model} $\subw_{\word}$ is the \kl{$\sig$-algebra} defined by:
    $\domain{\alg} = \wp(\Subw(\word))$,
    $\id^{\alg} = \set{\eps}$,
    $\emp^{\alg} = \emptyset$,
    and for all $\la[1], \la[2] \subseteq \Subw(\word)$,
    \begin{align*}
        \la[1] \compo^{\alg} \la[2] & = (\la[1] \compo \la[2]) \cap \Subw(\word),             &
        \la[1] \union^{\alg} \la[2] & = \la[1] \cup \la[2],                      \\
        \la[1]^{\kstar^{\alg}}      & = \la[1]^{\kstar} \cap \Subw(\word),                    &
        \la[1]^{\compl^{\alg}}      & = \Subw(\word) \setminus \la[1]. 
    \end{align*}
    We write $\SUBW_{\word}$ for the class of all \kl{valuations} of $\subw_{\word}$ and $\SUBW$.
    \lipicsEnd
\end{definition}
Let $\kappa_{\word}'$ be the map given by $L \mapsto L \cap \Subw(\word)$.
The map $\kappa_{\word}'$ forms an $\sig$-homomorphism from $\lang_{X}$ to $\subw_{\word}$.

We use $\ell_{1}, \ell_{2}, \dots$ as fixed pairwise distinct \kl{letters}.
We write $\SUBW$ for $\bigcup_{n \ge 0} \SUBW_{\ell_{1} \ell_{2} \dots \ell_{n}}$.
Let $\chi'$ be the map given by $L \mapsto \set{\ell_1 \dots \ell_{n} \mid a_1 \dots a_n \in L}$.
For each \kl{word} $a_1 \dots a_n$,
the map $\chi'$ forms an $\sig$-homomorphism from $\subw_{a_1 \dots a_n}$ to $\subw_{\ell_1 \dots \ell_n}$.
\begin{remark}
    The map $\chi'$ is defined, inspired by \kl{words-to-letters valuations} \cite{nakamuraWordstoLettersValuationsLanguage2023} where we consider ``\intro*\kl{letters-to-letters valuations}'' by restricting the mapping from \kl{words} to \kl{letters}.
\end{remark}
% Using this, we can embed the membership in $\LANG$ into that in $\SUBW$.
% \begin{proposition}\label{proposition: letter-to-letter valuations}
%     Let $X$ be a set and let $\val \in \LANG_{X}$.
%     Let $n \in \nat$ and let $a_0, \dots, a_{n-1}$ be \kl{letters} over $X$.
%     For all \kl{terms} $\term$ and $0 \le i \le j \le n$, $\ell_{i} \dots \ell_{j-1} \in \hat{\val}^{\tuple{a_0, \dots, a_{n-1}}}(\term)$ iff $a_{i} \dots a_{j-1} \in \hat{\val}(\term)$.
% \end{proposition}
% \begin{proof}
%     By easy induction on $\term$, cf.\ {\cite[Lem.\ 23]{nakamuraWordstoLettersValuationsLanguage2023}}\footnote{This equivalence holds because $a_i, \dots, a_{j-1}$ are \kl{letters} not \kl{words}. For \kl{words-to-letters valuations}, only the direction $\Rightarrow$ of \Cref{proposition: letter-to-letter valuations} holds {\cite[Lem.\ 23]{nakamuraWordstoLettersValuationsLanguage2023}}.} (see \Cref{section: proposition: letter-to-letter valuations}, for detailed proof).
%     %     We only write for Case $\term = \term[2] \compo \term[3]$ (see \Cref{section: proposition: letter-to-letter valuations} for each case).
%     %     Then we have
%     %     \begin{align*}
%     %          & \ell_{i} \dots \ell_{j-1} \in \hat{\val}^{\tuple{a_0, \dots, a_{n-1}}}(\term[2] \compo \term[3])                                                                                                                                                                       \\
%     %          & \hspace{-2.3em}\Longleftrightarrow \ell_{i} \dots \ell_{j-1} \in \hat{\val}^{\tuple{a_0, \dots, a_{n-1}}}(\term[2]) \compo \hat{\val}^{\tuple{a_0, \dots, a_{n-1}}}(\term[3]) \tag{$\ell_i \dots \ell_{j-1}$ is a subword of $\ell_{0} \dots \ell_{n-1}$}              \\
%     %          & \hspace{-2.3em}\Longleftrightarrow \bigvee_{i \le k \le j} (\ell_{i} \dots \ell_{k-1} \in \hat{\val}^{\tuple{a_0, \dots, a_{n-1}}}(\term[2]) \land \ell_{k} \dots \ell_{j-1} \in \hat{\val}^{\tuple{a_0, \dots, a_{n-1}}}(\term[3])) \tag{Def.\ of \kl{concatenation}} \\
%     %          & \hspace{-2.3em}\Longleftrightarrow \bigvee_{i \le k \le j} (a_{i} \dots a_{k-1} \in \hat{\val}(\term[2]) \land a_{k} \dots a_{j-1} \in \hat{\val}(\term[3])) \tag{IH}                                                                                                  \\
%     %          & \hspace{-2.3em}\Longleftrightarrow a_{i} \dots a_{j-1} \in \hat{\val}(\term[2]) \compo \hat{\val}(\term[3]). \tag*{(Def.\ of \kl{concatenation})\footnotemark \qedhere}
%     %     \end{align*} %
% \end{proof}
% Moreover, this transformation is a surjective map to $\SUBW$, as follows:
% \begin{proposition}\label{proposition: surjective}
%     Let $n \in \nat$ and let $\val' \in \SUBW_{n}$.
%     Let $\val \in \LANG_{\set{\ell_{0}, \dots, \ell_{n-1}}}$ be the \kl{valuation} defined by $\val(x) = \val'(x)$.
%     Then $\val^{\tuple{\ell_{0}, \dots, \ell_{n-1}}} = \val'$.
% \end{proposition}
% \begin{proof}
%     By \Cref{definition: word-to-letter valuations},
%     $\val^{\tuple{\ell_{0}, \dots, \ell_{n-1}}}(x) = \val(x) \cap \set{\ell_{i} \dots \ell_{j-1} \mid 0 \le i \le j \le n} = \val'(x)$.
% \end{proof}
% From them, we have that the \kl{quantifier-free theory} of $\LANG$ coincides with that of $\SUBW$.
% \begin{lemma}\label{lemma: LANG and SUBW}
%     For all \kl{quantifier-free formulas} $\fml$, we have: $\LANG \models \fml \ \Leftrightarrow\  \SUBW \models \fml$.
% \end{lemma}
% \begin{proof}
%     By $\val \models \term[1] = \term[2] \leftrightarrow (\term[1] \le \term[2] \land \term[2] \le \term[1])$ and $\val \models \term[1] \le \term[2] \leftrightarrow \term[1] \cap \term[2]^{-} \le \emp$ (which are set-theoretic equivalences) for $\val \in \LANG \cup \SUBW$, without loss of generality, we can assume that each \kl{equation} in $\fml$ is of the form $\term[3] \le \emp$.
%     By taking the conjunctive normal form, it suffices to show the claimed lemma when $\fml$ is of the form $(\bigvee_{k} \term[1]_{k} \le \emp) \lor (\bigvee_{k} \lnot \term[2]_{k} \le \emp)$.
%     By $\val \models (\term[1] \le \emp \land \term[2] \le \emp) \leftrightarrow \term[1] \cup \term[2] \le \emp$ (and $\val \models \mathsf{f} \leftrightarrow \top \le \emp$) for $\val \in \LANG \cup \SUBW$,
%     it suffices to show when $\fml$ is of the form $(\bigvee_{k = 1}^{m} \term[1]_{k} \le \emp) \lor (\lnot \term[2] \le \emp)$.
%     Then, we have:
%     $\SUBW \not\models (\bigvee_{k = 1}^{m} \term[1]_{k} \le \emp) \lor (\lnot \term[2] \le \emp)$
%     $\Rightarrow$ $\LANG \not\models (\bigvee_{k = 1}^{m} \term[1]_{k} \le \emp) \lor (\lnot \term[2] \le \emp)$ (by \Cref{proposition: surjective,proposition: letter-to-letter valuations})
%     $\Rightarrow$ $\exists n \in \nat, i_1, j_1, \dots, i_m, j_m \mbox{ s.t.\ $0 \le i_1 \le j_1 \le i_2 \le j_2 \le \dots \le i_m \le j_m \le n$}, \exists X, \val \in \LANG_{X}, a_{0}, \dots, a_{n-1} \in X,\ (\bigwedge_{k = 1}^{m}  a_{i_k} \dots a_{j_k-1} \in \hat{\val}(\term[1]_k)) \land \hat{\val}(\term[2]) = \emptyset$
%     $\Rightarrow$ $\SUBW \not\models (\bigvee_{k = 1}^{m} \term[1]_{k} \le \emp) \lor (\lnot \term[2] \le \emp)$ (by \Cref{proposition: letter-to-letter valuations} and $\val^{\tuple{a_0, \dots, a_{n-1}}} \in \SUBW_{n}$).
% \end{proof}

\begin{proposition}\label{proposition: subw and subwrel iso}
    For each $n \in \nat$, $\subw_{\ell_1 \dots \ell_{n}}$ is \kl{isomorphic} to $\rsub_n$.
\end{proposition}
\begin{proof}
    By the map $L \mapsto \set{\tuple{i,j} \mid 1 \le i \le j \le n+1 \land \ell_{i} \dots \ell_{j-1} \in L}$.
    This map can be also given by $\theta \circ \iota_{\set{\ell_1, \dots, \ell_{n}}}$.
    % As $f(\set{\id}) = \triangle_{\range{0, n}}$ (cf.\ \Cref{definition: relational bounded language model})
    % and $f(\set{\ell_{i} \dots \ell_{j-1}}) = \set{\tuple{i, j}}$ is singleton if $i < j$, this map is bijective.
    Fig.\ \ref{figure: proposition: subw and subwrel iso} presents illustrative instances of this map.
    %  (we display binary relations as graphs in a standard way).
\end{proof}
By \Cref{proposition: subw and subwrel iso}, the isomorphism closure of $\RSUB$ coincides with that of $\SUBW$.
Particularly, the \kl{quantifier-free theory} w.r.t.\ $\RSUB$ coincides with that of $\SUBW$.

\begin{figure}[ht]
    \begin{align*}
        \left\{\begin{aligned}
                   \eps
               \end{aligned}\right\}                                            & \qquad\mapsto\qquad \begin{tikzpicture}[baseline = -.5ex]
                                                                                                          \graph[grow right = 1.cm, branch down = 2.5ex, nodes={}]{
                                                                                                          {0/{$0$}[mysmallnode]}
                                                                                                          -!- {1/{$1$}[mysmallnode]}
                                                                                                          -!- {2/{$2$}[mysmallnode]}
                                                                                                          -!- {3/{$3$}[mysmallnode]}
                                                                                                          };
                                                                                                          \graph[use existing nodes, edges={color=black, pos = .5, earrow}, edge quotes={fill=white, inner sep=1pt,font= \scriptsize}]{
                                                                                                          {0,1,2,3} ->[out=-100, in=-80, looseness=15] {0,1,2,3};
                                                                                                          };
                                                                                                      \end{tikzpicture} \\[-1.ex]
        % \left\{\begin{aligned}
        %            \ell_{1}\ell_{2}
        %        \end{aligned}\right\}                                            & \qquad\mapsto\qquad \begin{tikzpicture}[baseline = -.5ex]
        %                                                                                                   \graph[grow right = 1.cm, branch down = 2.5ex, nodes={}]{
        %                                                                                                   {0/{$0$}[mysmallnode]}
        %                                                                                                   -!- {1/{$1$}[mysmallnode]}
        %                                                                                                   -!- {2/{$2$}[mysmallnode]}
        %                                                                                                   -!- {3/{$3$}[mysmallnode]}
        %                                                                                                   };
        %                                                                                                   \graph[use existing nodes, edges={color=black, pos = .5, earrow}, edge quotes={fill=white, inner sep=1pt,font= \scriptsize}]{
        %                                                                                                   1 ->[bend left = 20] 3;
        %                                                                                                   };
        %                                                                                               \end{tikzpicture} \\
        \left\{\begin{aligned}
                   \ell_{1}, \ell_{2}\ell_{3},\ \ell_{1}\ell_{2}, \ell_{3} 
               \end{aligned}\right\}                             & \qquad\mapsto\qquad \begin{tikzpicture}[baseline = -.5ex]
                                                                                           \graph[grow right = 1.cm, branch down = 2.5ex, nodes={}]{
                                                                                           {0/{$0$}[mysmallnode]}
                                                                                           -!- {1/{$1$}[mysmallnode]}
                                                                                           -!- {2/{$2$}[mysmallnode]}
                                                                                           -!- {3/{$3$}[mysmallnode]}
                                                                                           };
                                                                                           \graph[use existing nodes, edges={color=black, pos = .5, earrow}, edge quotes={fill=white, inner sep=1pt,font= \scriptsize}]{
                                                                                           0 -> 1;
                                                                                           1 ->[bend left = 20] 3;
                                                                                        %    {0,1,2,3} ->[out=-100, in=-80, looseness=15] {0,1,2,3};
          0 ->[bend right = 20] 2;
          2 -> 3;                                                                                 };
                                                                                       \end{tikzpicture}                \\[-1.ex]
        \left\{\ell_{i} \dots \ell_{j-1} \mid 1 \le i \le j \le 4 \right\} & \qquad\mapsto\qquad \begin{tikzpicture}[baseline = -.5ex]
a                                                                                                    \graph[grow right = 1.cm, branch down = 2.5ex, nodes={}]{
                                                                                                    {0/{$0$}[mysmallnode]}
                                                                                                    -!- {1/{$1$}[mysmallnode]}
                                                                                                    -!- {2/{$2$}[mysmallnode]}
                                                                                                    -!- {3/{$3$}[mysmallnode]}
                                                                                                    };
                                                                                                    \graph[use existing nodes, edges={color=black, pos = .5, earrow}, edge quotes={fill=white, inner sep=1pt,font= \scriptsize}]{
                                                                                                    0 -> 1;
                                                                                                    1 -> 2;
                                                                                                    2 -> 3;
                                                                                                    0 ->[bend left = 20] 2;
                                                                                                    1 ->[bend left = 20] 3;
                                                                                                    0 ->[bend left = 25] 3;
                                                                                                    {0,1,2,3} ->[out=-100, in=-80, looseness=15] {0,1,2,3};
                                                                                                    };
                                                                                                \end{tikzpicture}
    \end{align*}
    \caption{Illustrative instances of the bijective map of \Cref{proposition: subw and subwrel iso}, when $n = 3$.}
    \label{figure: proposition: subw and subwrel iso}
\end{figure}
Summarizing the above and \Cref{section: relation}, for each \kl{word} $\word_0 = a_1 \dots a_{n} \in X^*$,
we have considered the following $\sig$-homomorphisms:

\begin{tikzpicture}[baseline = -.5ex]
    \graph[grow right = 3.cm, branch down = 8ex, nodes={}]{
    {0l/{$\lang_{X}$}, 0r/{$\rlang_{X}$}}
    -!- {1l/{$\subw_{\word_0}$}[xshift = 0cm], 1r/{$\rlang_{X} \restriction \Subw(\word_0)$}[xshift = 0cm]}
    -!- {2l/{$\subw_{\ell_{1} \dots \ell_{n}}$}[xshift = 4cm], 2r/{$\rlang_{\set{\ell_{1}, \dots, \ell_{n}}} \restriction \Subw(\ell_{1} \dots \ell_{n})$}[xshift = 4cm], rsub/{$\rsub_{n}$}[xshift = 4cm]}
    };
    \graph[use existing nodes, edges={color=black, pos = .5, earrow}, edge quotes={fill=white, inner sep=1pt,font= \scriptsize}]{
        0l -> ["$\iota_{X}$" auto] 0r; 0l -> ["$\cong$"{auto, left, font = \tiny}] 0r;
        1l -> ["$\iota_{X}$" auto] 1r; 1l -> ["$\cong$"{auto, left, font = \tiny}] 1r;
        2l -> ["$\iota_{\set{\ell_{1}, \dots, \ell_{n}}}$" auto] 2r; 2l -> ["$\cong$"{auto, left, font = \tiny}] 2r;
    % {0,1,2,3} ->[out=-100, in=-80, looseness=15] {0,1,2,3};
        0l -> ["$\kappa'_{\word_0}$" auto] 1l; 0r -> ["$\kappa_{\word_0}$" auto] 1r;
        1l -> ["$\chi' \colon L \mapsto \set{\word \mid f(\word) \in L}$" auto] 2l;
        1r -> ["$R \mapsto \set{\tuple{\word, \word \word[2]} \mid \tuple{f(\word), f(\word \word[2])} \in R}$" auto, bend left = 10] 2r;
        2r -> ["$\theta$" auto] rsub; 2r -> ["$\cong$"{auto, left, font = \tiny}] rsub;
    };
\end{tikzpicture}\\
where $f$ is the monoid homomorphism given by $\ell_i \mapsto a_i$ for $i \in \range{1, n}$.
% \section{Proof of \Cref{proposition: letter-to-letter valuations}}\label{section: proposition: letter-to-letter valuations}
% (The following proof is almost the same as \cite[Lem.\ 23]{nakamuraWordstoLettersValuationsLanguage2023} except that we consider \kl{letters-to-letters valuations} instead of \kl{words-to-letters valuations}.)
% \begin{proof}[Proof of \Cref{proposition: letter-to-letter valuations}]
%        By easy induction on $\term$.

%        Case $\term = x$:
%        By the definition of $\val^{\tuple{a_0, \dots, a_{n-1}}}$,
%        we have
%        $\ell_{i} \dots \ell_{j-1} \in \hat{\val}^{\tuple{a_0, \dots, a_{n-1}}}(x) \Leftrightarrow a_{i}\dots a_{j-1} \in \hat{\val}(x)$.

%        Case $\term = \emp$:
%        By $\ell_{i} \dots \ell_{j-1} \in \hat{\val}^{\tuple{a_0, \dots, a_{n-1}}}(\emp) \Leftrightarrow \const{False} \Leftrightarrow  a_{i} \dots a_{j-1} \in \hat{\val}(\emp)$.

%        Case $\term = \id$:
%        By $\ell_{i} \dots \ell_{j-1} \in \hat{\val}^{\tuple{a_0, \dots, a_{n-1}}}(\id) \Leftrightarrow i = j \Leftrightarrow  a_{i} \dots a_{j-1} \in \hat{\val}(\id)$.

%        Case $\term = \term[2] \union \term[3]$:
%        By IH, we have
%        $\ell_{i} \dots \ell_{j-1} \in \hat{\val}^{\tuple{a_0, \dots, a_{n-1}}}(\term[2] \union \term[3])$
%        $\Leftrightarrow$ $\ell_{i} \dots \ell_{j-1} \in \hat{\val}^{\tuple{a_0, \dots, a_{n-1}}}(\term[2]) \lor \ell_{i} \dots \ell_{j-1} \in \hat{\val}^{\tuple{a_0, \dots, a_{n-1}}}(\term[3])$
%        $\Leftrightarrow$ $a_{i} \dots a_{j-1} \in \hat{\val}(\term[2]) \lor a_{i} \dots a_{j-1} \in \hat{\val}(\term[3])$
%        $\Leftrightarrow$ $a_{i} \dots a_{j-1} \in \hat{\val}(\term[2] \union \term[3])$.

%        Case $\term = \term[2]^{\compl}$:
%        By IH, we have
%        $\ell_{i} \dots \ell_{j-1} \in \hat{\val}^{\tuple{a_0, \dots, a_{n-1}}}(\term[2]^{\compl})$
%        $\Leftrightarrow$ $\ell_{i} \dots \ell_{j-1} \not\in \hat{\val}^{\tuple{a_0, \dots, a_{n-1}}}(\term[2])$
%        $\Leftrightarrow$ $a_{i} \dots a_{j-1} \not\in \hat{\val}(\term[2])$
%        $\Leftrightarrow$ $a_{i} \dots a_{j-1} \in \hat{\val}(\term[2]^{\compl})$.

%        Case $\term = \term[2] \compo \term[3]$:
%        We have
%        \begin{align*}
%                & \ell_{i} \dots \ell_{j-1} \in \hat{\val}^{\tuple{a_0, \dots, a_{n-1}}}(\term[2] \compo \term[3])                                                                                                                                                        \\
%                & \Leftrightarrow \ell_{i} \dots \ell_{j-1} \in \hat{\val}^{\tuple{a_0, \dots, a_{n-1}}}(\term[2]) \compo \hat{\val}^{\tuple{a_0, \dots, a_{n-1}}}(\term[3]) \tag{By the definition of $\hat{\val}^{\tuple{a_0, \dots, a_{n-1}}}$ and that $\ell_i \dots \ell_{j-1}$ is a subword of $\ell_{0} \dots \ell_{n-1}$}              \\
%                & \Leftrightarrow \bigvee_{i \le k \le j} (\ell_{i} \dots \ell_{k-1} \in \hat{\val}^{\tuple{a_0, \dots, a_{n-1}}}(\term[2]) \land \ell_{k} \dots \ell_{j-1} \in \hat{\val}^{\tuple{a_0, \dots, a_{n-1}}}(\term[3])) \tag{By \kl{concatenation}}               \\
%                & \Leftrightarrow \bigvee_{i \le k \le j} (a_{i} \dots a_{k-1} \in \hat{\val}(\term[2]) \land a_{k} \dots a_{j-1} \in \hat{\val}(\term[3])) \tag{IH}                                                                                                  \\
%                & \Leftrightarrow a_{i} \dots a_{j-1} \in \hat{\val}(\term[2] \compo \term[3]). \tag{By \kl{concatenation}}
%        \end{align*}
%        (Note that the last equivalence holds thanks to that $a_{i}, \dots, a_{j-1}$ are \kl{letters} not \kl{words}, cf.\ \cite[Lem.\ 23]{nakamuraWordstoLettersValuationsLanguage2023}.)

%        Case $\term = \term[2]^{\kstar}$:
%        We have
%        \begin{align*}
%                & \ell_{i} \dots \ell_{j-1} \in \hat{\val}^{\tuple{a_0, \dots, a_{n-1}}}(\term[2]^{\kstar})                                                                                                                                       \\
%                & \Leftrightarrow \ell_{i} \dots \ell_{j-1} \in \hat{\val}^{\tuple{a_0, \dots, a_{n-1}}}(\term[2])^{\kstar}  \tag{By the definition of $\hat{\val}^{\tuple{a_0, \dots, a_{n-1}}}$ and that $\ell_i \dots \ell_{j-1}$ is a \kl{subword} of $\ell_{0} \dots \ell_{n-1}$}                                      \\
%                & \Leftrightarrow \exists m \in \nat, \bigvee_{i = k_0 \le k_1 \le \dots \le k_m = j} \bigwedge_{l = 1}^{m} (\ell_{k_{(l-1)}} \dots \ell_{k_{l}-1} \in \hat{\val}^{\tuple{a_0, \dots, a_{n-1}}}(\term[2]))  \tag{By \kl{Kleene-star}} \\
%                & \Leftrightarrow \exists m \in \nat, \bigvee_{i = k_0 \le k_1 \le \dots \le k_m = j} \bigwedge_{l = 1}^{m} (a_{k_{(l-1)}} \dots a_{k_{l}-1} \in \hat{\val}(\term[2])) \tag{IH}                                                 \\
%                & \Leftrightarrow a_{i} \dots a_{j-1} \in \hat{\val}(\term[2]^{\kstar}).  \tag*{(By \kl{Kleene-star}) \qedhere}
%        \end{align*}
% \end{proof}



\section{Explicit proof of \Cref{lemma: CFG and LANG}}\label{section: proofs in undecidability}
(The following proof is almost the same as \cite[Section VI]{nakamuraExistentialCalculiRelations2023} except that we consider $\RSUB$ instead of $\REL$.)

For a \kl{CFG} $\CFG = \tuple{X, \mathcal{R}, s}$ and a \kl{word} $\word = a_1 \dots a_n$,
let $\val_{\CFG, \word} \in \RSUB_{n}$ be the \kl{valuation} defined by:
\begin{itemize}
    \item $\val_{\CFG, \word}(x) = \set{\tuple{i-1, i} \mid i \in \range{1, n} \land x = a_i}$ for $x \in \CSIG$;
    \item $\set{\val_{\CFG, \word}(x)}_{x \in X}$ are the minimal sets such that
          for all $x \in X$, $\word[2] \in (X \cup \CSIG)^*$, and $i, j \in \range{0, n}$, if $x \leftarrow \word[2] \in \mathcal{R}$, then
          $\begin{prooftree}
                  \hypo{\tuple{i,j} \in \hat{\val}_{\CFG, \word}(\word[2])}
                  \infer1{\tuple{i,j} \in \val_{\CFG, \word}(x)}
              \end{prooftree}$.
\end{itemize}
For instance, if $A = \set{(,)}$,
$\CFG = \tuple{\set{s}, \set{s \leftarrow (s) s,\  s \leftarrow \const{I}}, s}$ (i.e., $\ljump{\CFG}$ is the Dyck-1 language), and $w = (()())$, then $\val_{\CFG, \word}$ is of the following form:
\[\val_{\CFG, \word} \quad=\quad \begin{tikzpicture}[baseline = -.5ex]
        \graph[grow right = 1.2cm, branch down = 6ex, nodes={mysmallnode, font = \scriptsize}]{
        {0/{$0$}[draw, circle]}
        -!-{1/{$1$}[draw, circle]}
        -!-{2/{$2$}[draw, circle]}
        -!-{3/{$3$}[draw, circle]}
        -!-{4/{$4$}[draw, circle]}
        -!-{5/{$5$}[draw, circle]}
        -!-{6/{$6$}[draw, circle]}
        };
        \path (0) edge [draw = white, opacity = 0] node[pos= 0.5, inner sep = 1.5pt, font = \scriptsize, opacity = 1](a01){$($}(1);
        \path (1) edge [draw = white, opacity = 0] node[pos= 0.5, inner sep = 1.5pt, font = \scriptsize, opacity = 1](a12){$($}(2);
        \path (2) edge [draw = white, opacity = 0] node[pos= 0.5, inner sep = .5pt, font = \scriptsize, opacity = 1](a23){$)$}(3);
        \path (3) edge [draw = white, opacity = 0] node[pos= 0.5, inner sep = .5pt, font = \scriptsize, opacity = 1](a34){$($}(4);
        \path (4) edge [draw = white, opacity = 0] node[pos= 0.5, inner sep = .5pt, font = \scriptsize, opacity = 1](a45){$)$}(5);
        \path (5) edge [draw = white, opacity = 0] node[pos= 0.5, inner sep = .5pt, font = \scriptsize, opacity = 1](a56){$)$}(6);
        \node(a00)[above = 1.5ex of 0,  inner sep = 1.5pt, font = \scriptsize,]{$s$};
        \node(a11)[above = 1.5ex of 1,  inner sep = 1.5pt, font = \scriptsize,]{$s$};
        \node(a22)[above = 1.5ex of 2,  inner sep = 1.5pt, font = \scriptsize,]{$s$};
        \node(a33)[above = 1.5ex of 3,  inner sep = 1.5pt, font = \scriptsize,]{$s$};
        \node(a44)[above = 1.5ex of 4,  inner sep = 1.5pt, font = \scriptsize,]{$s$};
        \node(a55)[above = 1.5ex of 5,  inner sep = 1.5pt, font = \scriptsize,]{$s$};
        \node(a66)[above = 1.5ex of 6,  inner sep = 1.5pt, font = \scriptsize,]{$s$};
        % \path (1) edge [draw = white, opacity = 0, bend left = 45] node[pos= 0.5, inner sep = 1.5pt, font = \scriptsize, opacity = 1](I13){$\const{I}$}(3);
        \graph[use existing nodes, edges={color=black, pos = .5, earrow}, edge quotes={fill=white, inner sep=1pt,font= \scriptsize}]{
        0 -- a01 -> 1;
        1 -- a12 -> 2;
        2 -- a23 -> 3;
        3 -- a34 -> 4;
        4 -- a45 -> 5;
        5 -- a56 -> 6;
        0 --[bend left = 15] a00 ->[bend left = 15] 0;
        1 --[bend left = 15] a11 ->[bend left = 15] 1;
        2 --[bend left = 15] a22 ->[bend left = 15] 2;
        3 --[bend left = 15] a33 ->[bend left = 15] 3;
        4 --[bend left = 15] a44 ->[bend left = 15] 4;
        5 --[bend left = 15] a55 ->[bend left = 15] 5;
        6 --[bend left = 15] a66 ->[bend left = 15] 6;
        1 ->["$s$", font=\scriptsize, bend right = 20] 3;
        3 ->["$s$", font=\scriptsize, bend right = 20] 5;
        1 ->["$s$", font=\scriptsize, bend right = 20] 5;
        0 ->["$s$", font=\scriptsize, bend right = 20] 6;
        };
    \end{tikzpicture}\]
By construction, we have the following:
\begin{lemma}[cf.\ {\cite[Lem.\ 47]{nakamuraExistentialCalculiRelations2023}}]\label{lemma: CFG}
    Let $\CFG = \tuple{X, \mathcal{R}, s}$ be a \kl{CFG} over a finite set $\CSIG$.
    Let $\algclass$ be such that $\set{\val_{\CFG, \word} \mid \word \in \CSIG^*} \subseteq \algclass \subseteq \GREL$.
    For all $y \in X$ and $\word[2] \in \CSIG^*$, we have:
    \[y \vdash_{\CFG} \word[2] \quad\Leftrightarrow\quad \algclass \models (\bigwedge_{(x \leftarrow w) \in \mathcal{R}} w \le x) \to \word[2] \le y.\]
\end{lemma}
\begin{proof}
    ($\Rightarrow$):
    By induction on the derivation tree of $\vdash_{\CFG}$.
    Let us consider \[\begin{prooftree}
                  \hypo{x_1 \vdash_{\CFG} \word[2]_1}
                  \hypo{\dots}
                  \hypo{x_m \vdash_{\CFG} \word[2]_m}
                  \infer3{y \vdash_{\CFG} \word[1]_0 \word[2]_1 \word[1]_1 \dots \word[2]_m \word[1]_m}
           \end{prooftree} \mbox{ where $y \leftarrow \word[1]_0 x_1 \word[1]_1 \dots x_m \word[1]_m \in \mathcal{R}$.}\]
    Then, w.r.t.\ ``$\val_{\CFG, \word} \models \bl$'',
    we have:
    \begin{align*}
        \word[1]_0 \word[2]_1 \word[1]_1 \dots \word[2]_m \word[1]_m & \le \word[1]_0 x_1 \word[1]_1 \dots x_m \word[1]_m \tag{By IH, $\val_{\CFG, \word} \models \word[2]_k \le x_k$}\\
        & \le y. \tag{Because $\val_{\CFG, \word} \models \bigwedge_{(x \leftarrow w) \in \mathcal{R}} w \le x$}
    \end{align*}

    ($\Leftarrow$):
    Let $\word[2] = a_0 \dots a_{n-1}$.
    By construction, we have $\val_{\CFG, \word[2]} \models \bigwedge_{(x \leftarrow w) \in \mathcal{R}} w \le x$ and $\tuple{0, n} \in \hat{\val}_{\CFG, \word[2]}(\word[2])$, so, we have $\tuple{0, n} \in \val_{\CFG, \word[2]}(y)$.
    Then, it suffices to show the following:
    $\forall i, j \in \range{0, n}, \forall y \in X$, $\tuple{i, j} \in \val_{\CFG, \word[2]}(y) \Longrightarrow y \vdash_{\CFG} a_i \dots a_{j-1}$.
    We show this by induction on the derivation tree induced from the definition of $\val_{\CFG, \word[2]}$.
    Let us consider
    \[\begin{prooftree}
        \hypo{\tuple{i, j} \in \hat{\val}_{\CFG, \word[2]}(\word[1]_0 x_1 \word[1]_1 \dots x_m \word[1]_m)}
        \infer1{\tuple{i, j} \in \val_{\CFG, \word[2]}(y)}
    \end{prooftree} \mbox{ where $y \leftarrow \word[1]_0 x_1 \word[1]_1 \dots x_m \word[1]_m \in \mathcal{R}$.}\]
    Let $i = i_0 \le j_0 \le \dots \le i_{m} \le j_{m} = j$ be s.t.\
    $\tuple{i_k, j_k} \in \hat{\val}_{\CFG, \word[2]}(\word[1]_k)$ and
    $\tuple{j_{k-1}, i_k} \in \val_{\CFG, \word[2]}(x_k)$.
    By the form of $\val_{\CFG, \word[2]}$, we have $\word[1]_k = a_{i_k} \dots a_{j_k-1}$.
    By IH, we have $x_k \vdash_{\CFG} a_{j_{k-1}} \dots a_{i_k-1}$.
    Combining them yields $y \vdash_{\CFG} a_{i} \dots a_{j-1}$.
\end{proof}
\begin{lemma}[cf.\ {\cite[Lem.\  48]{nakamuraExistentialCalculiRelations2023}}]\label{lemma: CFG encoding lang}
    Let $\CFG = \tuple{X, \mathcal{R}, s}$ be a \kl{CFG} over a finite set $\CSIG = \set{a_1, \dots, a_n}$.
    Let $\algclass$ be such that $\set{\val_{\CFG, \word} \mid \word \in \CSIG^*} \subseteq \algclass \subseteq \GREL$.
    Then we have:
    \[\ljump{\CFG} = \CSIG^* \quad\Leftrightarrow\quad\algclass \models  (\bigwedge_{(x \leftarrow w) \in \mathcal{R}} w \le x) \to (\sum_{i = 1}^{n} a_i)^* \le s.\]
    % Here, $\CSIG^*$ denotes the term $(\sum_{i = 1}^{n} a_i)^*$ in the right-hand side, where $\CSIG = \set{a_1, \dots, a_n}$.
\end{lemma}
\begin{proof}
    ($\Rightarrow$):
    By \Cref{lemma: CFG} with that $\GREL$ is $*$-continuous (i.e., $\displaystyle x y^* z = \sup_{n \ge 0} x y^{n} z$).
    ($\Leftarrow$):
    By \Cref{lemma: CFG} with $\algclass \models \word[2] \le (\sum_{i = 1}^{n} a_i)^*$ for all $\word[2] \in \CSIG^*$.
\end{proof}
\begin{proof}[Proof of \Cref{lemma: CFG and LANG}]
    By \Cref{lemma: CFG encoding lang} (since $\set{\val_{\CFG, \word} \mid \word \in \CSIG^*} \subseteq \RSUB$).
\end{proof}

\begin{remark}
    \Cref{lemma: CFG encoding lang} holds for $\LANG$ by viewing hypotheses as fixed-point equations.
    We use $\RSUB$ for using Hore hypotheses elimination \Cref{corollary: ORD lr Hoare} and for showing \Cref{lemma: CFG and LANG star-free}.
\end{remark}

\section{Proof of \Cref{lemma: CFG and LANG star-free}}\label{section: lemma: CFG and LANG star-free}
\begin{proof}
    Because $\val_{\CFG, \word} \models \com{\id} = (\sum_{i = 1}^{n} a_i) \top$, we have $\set{\val_{\CFG, \word} \mid \word \in \CSIG^*} \subseteq \set{\val \in \RSUB \mid \val \models \com{\id} = (\sum_{i = 1}^{n} a_i) \top}$.
    Thus, by \Cref{lemma: CFG encoding lang} with applying \Cref{proposition: star to top} to $(\sum_{i = 1}^{n} a_i)^*$,
    we have:
    \[\ljump{\CFG} = \CSIG^* \quad\Leftrightarrow\quad \RSUB \models (\com{\id} = (\sum_{i = 1}^{n} a_i) \top \land \bigwedge_{(x \leftarrow w) \in \mathcal{R}} w \le x) \to \top \le s. \qedhere\] 
\end{proof}

\section{Trakhtenbrot's theorem from the universality problem of CFGs (\Cref{remark: Trakhtenbrot})}\label{section: remark: Trakhtenbrot}
In this section, we note that Trakhtenbrot's theorem \cite{trakhtenbrotImpossibilityAlgorithmDecision1950} can be derived as an analog of \Cref{section: undecidability}.
Let $\REL^{\mathrm{fin}}$ be the class of all \kl{valuations} of all \kl{relational models} on \emph{finite} sets.

Let $z$ be a fresh \kl{variable} and let $\fml_0$ be the following \kl{formula}:
\begin{align*}
    \fml_0 &\quad\defeq\quad (z \le \com{\id}) \land (z z \le z) \land (\top \le \top ((\com{z} \cap \com{\id}) \top)^{-})  \land (z \le zz \union \sum_{i = 1}^{n} a_i).
\end{align*}
This \kl{formula} is defined based on \cite[p.\ 30]{ebbinghausFiniteModelTheory1995};
the relation $z$ is irreflexive (by $z \le \com{\id}$) and transitive (by $z z \le z$) with a minimal element (by $\top \le \top ((\com{z} \cap \com{\id}) \top)^{-}$),
and that $\tuple{x, y} \in \hat{\val}(\sum_{i = 1}^{n} a_i)$ holds for any immediate successor $y$ of $x$ (by $z \le zz \union \sum_{i = 1}^{n} a_i$ in finite models).
Note that by construction, $\REL^{\mathrm{fin}} \models \fml_0 \to z \union \id \le (\sum_{i = 1}^{n} a_i)^*$.

\begin{lemma}\label{lemma: CFG encoding lang Trakhten}
    Let $\CFG = \tuple{X, \mathcal{R}, s}$ be a \kl{CFG} over a finite set $\CSIG = \set{a_1, \dots, a_n}$.
    Then we have:
    \[\ljump{\CFG} = \CSIG^* \quad\Leftrightarrow\quad \REL^{\mathrm{fin}} \models (\fml_0 \land \bigwedge_{(x \leftarrow w) \in \mathcal{R}} w \le x) \to z \union \id \le s.\]
\end{lemma}
\begin{proof}
    ($\Rightarrow$):
    By \Cref{lemma: CFG encoding lang} with $\REL^{\mathrm{fin}} \models \fml_0 \to z \union \id \le (\sum_{i = 1}^{n} a_i)^*$.
    ($\Leftarrow$):
    Let $\word[2] \in \CSIG^*$ be arbitrary.
    Let $\val_{\CFG, \word[2]}' \in \REL^{\mathrm{fin}}$ be the \kl{valuation} on the \kl{relational model} on $\range{0, n}$ such that $\val_{\CFG, \word[2]}'(x) = \val_{\CFG, \word[2]}(x)$ for $x \in \CSIG \cup X$ and $\val_{\CFG, \word[2]}'(z) = \hat{\val}_{\CFG, \word[2]}(\com{\id})$. 
    Then, one can see that $\val_{\CFG, \word[2]}' \models \fml_0$,
    $\val_{\CFG, \word[2]}' \models \bigwedge_{(x \leftarrow w) \in \mathcal{R}} w \le x$, thus we have $\val_{\CFG, \word[2]}' \models z \union \id \le s$.
    By $\val_{\CFG, \word[2]}' \models z \union \id = (\sum_{i = 1}^{n} a_i)^* \ge \word[2]$, we have $\val_{\CFG, \word[2]}' \models \word[2] \le s$.
    By $\hat{\val}_{\CFG, \word[2]}'(\word[2]') = \hat{\val}_{\CFG, \word[2]}(\word[2]')$ for $\word[2]' \in (\CSIG \cup X)^*$, we also have $\val_{\CFG, \word[2]} \models \word[2] \le s$.
    Thus by \Cref{lemma: CFG}, we have $\word[2] \in \ljump{\CFG}$.
    Hence, $\ljump{\CFG} = \CSIG^*$.
\end{proof}

By \Cref{lemma: CFG encoding lang Trakhten}, we have that the \kl{quantifier-free theory} of \kl{terms} over $\set{\union, \compo, \bl^{\compl}, \id}$ is $\Pi^{0}_{1}$-complete w.r.t.\ $\REL^{\mathrm{fin}}$.
Thus, (e.g.,  by the standard translation \cite{tarskiCalculusRelations1941} from \kl{quantifier-free formulas} of the calculus of relations into the formulas of first-order logic),
we have that the finite validity problem of first-order logic is $\Pi^{0}_{1}$-complete (Trakhtenbrot's theorem \cite{trakhtenbrotImpossibilityAlgorithmDecision1950}).



\section{Explicit Proof of \Cref{theorem: graph characterization}}\label{section: theorem: graph characterization}
Let $\val \in \GREL$.
For a \kl{graph} $\graph[2]$ and a \kl{graph language} $\glang$, we write
\begin{align*}
    \hat{\val}(\graph[2]) & \quad\defeq\quad \set{\tuple{x, y} \mid \graph[2] \homo \const{\graph}(\val, x, y)}, & \hat{\val}(\glang) & \quad\defeq\quad \bigcup_{\graph[2] \in \glang} \hat{\val}(\graph[2]).
\end{align*}
\begin{lemma}\label{lemma: glang val}
    Let $\val \in \GREL$.
    For all $\KAtermclass_{\set{\com{x}, \com{\id}, \top, \cap}}$ \kl{terms} $\term$,
    we have $\hat{\val}(\term) = \hat{\val}(\glang(\term))$.
\end{lemma}
\begin{proof}[Proof (cf.\ {\cite[Prop.\  11]{nakamuraExistentialCalculiRelations2023}})]
    By easy induction on $\term$.
    See \Cref{section: lemma: glang val} for details.
\end{proof}

\begin{lemma}\label{lemma: QS}
    Let $\algclass \subseteq \GREL$, $\val \in \algclass$, and $\graph[2]$ be a \kl{graph}.
    Then $\hat{\val}(\FillS_{\algclass}(\graph[2])^{\FillQ}) = \hat{\val}(\graph[2])$.
\end{lemma}
\begin{proof}
    ($\subseteq$):
    Because, for any $\graph[3] \in \FillS_{\algclass}(\graph[2])$, 
    we have $\graph[3] \homo \graph[3]^{\FillQ}$ by $\graph[2] \homo \graph[3]$ and $\graph[3] \homo \graph[3]^{\FillQ}$.
    % Hence $\hat{\val}(\FillS_{\algclass}(\graph[2])^{\FillQ}) \subseteq \hat{\val}(\graph[2])$.
    ($\supseteq$):
    Let $h \colon \graph[2] \homo \const{\graph}(\val, x, y)$.
    Then $\FillS(h) \homo \const{\graph}(\val, x, y)$.
    Because $\id^{\const{\graph}(\val, x, y)}$ is the identity relation, $\FillS(h)^{\FillQ} \homo \const{\graph}(\val, x, y)$.
    Hence $\hat{\val}(\FillS_{\algclass}(\graph[2])^{\FillQ}) \supseteq \hat{\val}(\graph[2])$.
\end{proof}


% \begin{proposition}\label{proposition: glang decomposition}
%     For all $\KAtermclass_{\set{\com{x}, \com{\id}, \cap}}$ terms $\term[1], \term[2]$, we have
%     \[\RSUB \models \term[1] \le \term[2] \quad\Leftrightarrow\quad \forall \graph \in \glang(\term[1]),\  \RSUB \models \graph[1] \le \term[2].\]
% \end{proposition}
% \begin{proof}
%     For all $\val \in \VAL_{\RSUB}$, we have
%     \begin{align*}
%         \hat{\val}(\term[1]) \subseteq \hat{\val}(\term[2]) & \quad\Leftrightarrow\quad \hat{\val}(\glang(\term[1])) \subseteq \hat{\val}(\term[2])             \tag{\Cref{proposition: glang val}} \\
%                                                             & \quad\Leftrightarrow\quad \forall \graph \in \glang(\term[1]),\  \hat{\val}(\graph) \subseteq \hat{\val}(\term[2]).
%     \end{align*}
%     Hence, this completes the proof.
% \end{proof}

\begin{proposition}\label{proposition: QS and ORD}
    Let $h \colon \graph[2] \homo \const{\graph}(\val, x, y)$.
    Then the \kl{graph} $\FillS(h)^{\FillQ}$ is \kl{isomorphic} to the \kl{induced subgraph} of $\const{\graph}(\val, x, y)$ on the range of $h$.
\end{proposition}
\begin{proof}
    Easy, by construction.
    Note that since $\id^{\const{\graph}(\val, x, y)}$ is the identity relation, $\id^{\FillS(h)}$ is an equivalence relation (see, e.g., \Cref{example: edge-saturation}).
\end{proof}


\begin{proof}[Proof of \Cref{theorem: graph characterization} (cf.\ {\cite[Thm.\ 18]{nakamuraExistentialCalculiRelations2023}})]
    We have:
    \begin{align*}
         & \algclass \models \term[1] \le \term[2] \ \Leftrightarrow\   \forall \val \in \algclass, \hat{\val}(\term[1]) \subseteq \hat{\val}(\term[2])                                                                                          \\
         & \Leftrightarrow \forall \val \in \algclass, \hat{\val}(\FillS_{\algclass}(\glang(\term[1]))^{\FillQ}) \subseteq \hat{\val}(\glang(\term[2])) \tag{\Cref{lemma: glang val,lemma: QS}}                                                                                                                                                      \\
         & \Leftrightarrow \forall \graph \in \FillS_{\algclass}(\glang(\term[1]))^{\FillQ}, \forall \graph[3] \in \allgraph_{\algclass},\  (\graph \homo \graph[3]) \mbox{ implies } (\exists \graph[2] \in \glang(\term[2]), \ \graph[2] \homo \graph[3])                                                \tag{Def.\ of $\hat{\val}$}         \\
         & \Leftrightarrow \forall \graph \in \FillS_{\algclass}(\glang(\term[1]))^{\FillQ}, \exists \graph[2] \in \glang(\term[2]), \ \graph[2] \homo \graph[1]. \tag{$\heartsuit$}
    \end{align*}
    Here, for ($\heartsuit$),
    ($\Leftarrow$):
    Let $\graph[2] \in \glang(\term[2])$ be s.t.\  $\graph[2] \homo \graph[1]$.
    Then for all $\graph[3]$ s.t.\ $\graph[1] \homo \graph[3]$,
    we have $\graph[2] \homo \graph[3]$ by transitivity of $\homo$.
    ($\Rightarrow$):
    By \Cref{proposition: QS and ORD} and that $\allgraph_{\algclass}$ is \kl{induced subgraph-closed},
    we have $\FillS_{\algclass}(\glang(\term[1]))^{\FillQ} \subseteq \allgraph_{\algclass}$.
    Thus by letting $\graph[3] = \graph$, this completes the proof.
\end{proof}


\subsection{Proof of \Cref{lemma: glang val}}\label{section: lemma: glang val}
(The following proof is almost the same as \cite[Prop.\ 8, 11]{nakamuraExistentialCalculiRelations2023} except that we consider $\GREL$ instead of $\REL$.
Similar arguments can also be found, e.g., in \cite{andrekaEquationalTheoryUnionfree1995,brunetPetriAutomata2017,pousPositiveCalculusRelations2018}.)

We use the two notations for \kl{graphs}, \kl{series-composition} ($\compo$) and \kl{parallel-composition} ($\cap$):
\begin{align*}
       G \compo H
        & \quad\defeq\quad \begin{tikzpicture}[baseline = -.5ex]
                        \graph[grow right = 1.cm, branch down = 2.5ex, nodes={mynode, font = \scriptsize}]{
                        {s1/{}[draw, circle]}
                        -!- {c/{}[draw, circle]}
                        -!- {t1/{}[draw, circle]}
                        };
                        \node[left = 4pt of s1](s1l){} edge[earrow, ->] (s1);
                        \node[right = 4pt of t1](t1l){}; \path (t1) edge[earrow, ->] (t1l);
                        \graph[use existing nodes, edges={color=black, pos = .5, earrow}, edge quotes={fill=white, inner sep=1pt,font= \scriptsize}]{
                        s1 ->["$G$"] c;
                        c ->["$H$"] t1;
                        };
                 \end{tikzpicture},
        & G \cap H
        & \quad\defeq\quad \begin{tikzpicture}[baseline = -.5ex]
                        \graph[grow right = 1.cm, branch down = 2.5ex, nodes={mynode, font = \scriptsize}]{
                        {s1/{}[draw, circle]}
                        -!- {t1/{}[draw, circle]}
                        };
                        \node[left = 4pt of s1](s1l){} edge[earrow, ->] (s1);
                        \node[right = 4pt of t1](t1l){}; \path (t1) edge[earrow, ->] (t1l);
                        \graph[use existing nodes, edges={color=black, pos = .5, earrow}, edge quotes={fill=white, inner sep=1pt,font= \scriptsize}]{
                        s1 ->["$G$", bend left = 25] t1;
                        s1 ->["$H$", bend right = 25] t1;
                        };
                 \end{tikzpicture}.
\end{align*}

\begin{proposition}\label{lemma: glang val op}
       Let $\val \in \GREL$ and $\graph[1], \graph[2]$ be \kl{graphs}.
       \begin{align*}
              \label{lemma: glang val cap} \hat{\val}(\graph[1] \cap \graph[2])     & = \hat{\val}(\graph[1]) \cap \hat{\val}(\graph[2]) \tag{\Cref{lemma: glang val op}$\cap$}     \\
              \label{lemma: glang val compo} \hat{\val}(\graph[1] \compo \graph[2]) & = \hat{\val}(\graph[1]) \rcompo \hat{\val}(\graph[2]) \tag{\Cref{lemma: glang val op}$\compo$}
       \end{align*}
\end{proposition}
\begin{proof}
       (\ref{lemma: glang val cap}):
       It suffices to prove that for every $x, y$,
       \[\exists f, f \colon (\graph[1] \cap \graph[2]) \homo \const{\graph}(\val, x, y) \;\Leftrightarrow\; \exists f_{\graph[1]}, f_{\graph[2]},\  f_{\graph[1]} \colon \graph[1] \homo \const{\graph}(\val, x, y) \land f_{\graph[2]} \colon \graph[2] \homo \const{\graph}(\val, x, y).\]
       ($\Rightarrow$):
       By letting $f_{\graph[1]} \defeq \set{\tuple{x', f(x')} \mid x' \in |\graph[1]|}$ and $f_{\graph[2]} \defeq \set{\tuple{x', f(x')} \mid x' \in |H|}$.
       ($\Leftarrow$):
       By letting $f \defeq f_{\graph[1]} \cup f_{\graph[2]}$.
       Note that $f_{\graph[1]}(\src^{\graph[1]}) = x = f_{\graph[2]}(\src^{\graph[2]})$ and $f_{\graph[1]}(\tgt^{\graph[1]}) = y = f_{\graph[2]}(\tgt^{\graph[2]})$; so $f$ is indeed a map.

       (\ref{lemma: glang val compo}):
       It suffices to prove that for every $x, y$,
       \[\exists f, f \colon (\graph[1] \compo \graph[2]) \homo \const{\graph}(\val, x, y) \;\Leftrightarrow\; \exists z, \exists f_{\graph[1]}, f_{\graph[2]},\  f_{\graph[1]} \colon \graph[1] \homo \const{\graph}(\val, x, z) \land f_{\graph[2]} \colon \graph[2] \homo \const{\graph}(\val, z, y).\]
       ($\Rightarrow$):
       By letting $z \defeq f(\tgt^{\graph[1]})$, $f_{\graph[1]} \defeq \set{\tuple{x', f(x')} \mid x' \in |\graph[1]|}$, and $f_{\graph[2]} \defeq \set{\tuple{x', f(x')} \mid x' \in |\graph[2]|}$.
       ($\Leftarrow$):
       By letting $f \defeq f_{\graph[1]} \cup f_{\graph[2]}$.
       Note that $f_{\graph[1]}(\tgt^{\graph[1]}) = z = f_{\graph[2]}(\src^{\graph[2]})$; so $f$ is indeed a map.
\end{proof}

\begin{proof}[Proof of \Cref{lemma: glang val}]
       By easy induction on $\term$.

       Case $\term = x$ where $x \in \tilde{\vsig} = \set{a, \com{a} \mid a \in \vsig} \cup \set{\com{\id}, \top}$:
       For every $\tuple{x, y} \in \hat{\val}(\top)$, we have
       \begin{align*}
              \tuple{x, y} \in \hat{\val}(a) & \quad\Leftrightarrow\quad (\begin{tikzpicture}[baseline = -.5ex]
                                                                     \graph[grow right = 1.2cm, branch down = 6ex, nodes={mynode}]{
                                                                     {0/{}[draw, circle]}-!-{1/{}[draw, circle]}
                                                                     };
                                                                     \node[left = .5em of 0](l){};
                                                                     \node[right = .5em of 1](r){};
                                                                     \path (0) edge [draw = white, opacity = 0] node[pos= 0.5, inner sep = 1.5pt, font = \scriptsize, opacity = 1](a1){$a$}(1);
                                                                     \graph[use existing nodes, edges={color=black, pos = .5, earrow}, edge quotes={fill=white, inner sep=1pt,font= \scriptsize}]{
                                                                            0 -- a1 -> 1;
                                                                            l -> 0; 1 -> r;
                                                                     };
                                                              \end{tikzpicture} \homo \const{\graph}(\val, x, y))                             \\
                                             & \quad\Leftrightarrow\quad \tuple{x, y} \in \hat{\val}(\begin{tikzpicture}[baseline = -.5ex]
                                                                                                 \graph[grow right = 1.2cm, branch down = 6ex, nodes={mynode}]{
                                                                                                 {0/{}[draw, circle]}-!-{1/{}[draw, circle]}
                                                                                                 };
                                                                                                 \node[left = .5em of 0](l){};
                                                                                                 \node[right = .5em of 1](r){};
                                                                                                 \path (0) edge [draw = white, opacity = 0] node[pos= 0.5, inner sep = 1.5pt, font = \scriptsize, opacity = 1](a1){$a$}(1);
                                                                                                 \graph[use existing nodes, edges={color=black, pos = .5, earrow}, edge quotes={fill=white, inner sep=1pt,font= \scriptsize}]{
                                                                                                        0 -- a1 -> 1;
                                                                                                        l -> 0; 1 -> r;
                                                                                                 };
                                                                                          \end{tikzpicture}) \tag{Def.\ of $\hat{\val}$} \\
                                             & \quad\Leftrightarrow\quad \tuple{x, y} \in \hat{\val}(\glang(a)). \tag{Def.\ of $\glang$}
       \end{align*}

       Case $\term = \id$:
       For every $\tuple{x, y} \in \hat{\val}(\top)$, we have
       \begin{align*}
              \tuple{x, y} \in \hat{\val}(\id) & \quad\Leftrightarrow\quad (\begin{tikzpicture}[baseline = -.5ex]
                                                                        \graph[grow right = 1.2cm, branch down = 6ex, nodes={mynode}]{
                                                                        {0/{}[draw, circle]}-!-{1/{}[draw, circle]}
                                                                        };
                                                                        \node[left = .5em of 0](l){};
                                                                        \node[right = .5em of 1](r){};
                                                                        \path (0) edge [draw = white, opacity = 0] node[pos= 0.5, inner sep = 1.5pt, font = \scriptsize, opacity = 1](a1){$\id$}(1);
                                                                        \graph[use existing nodes, edges={color=black, pos = .5, earrow}, edge quotes={fill=white, inner sep=1pt,font= \scriptsize}]{
                                                                               0 -- a1 -> 1;
                                                                               l -> 0; 1 -> r;
                                                                        };
                                                                 \end{tikzpicture} \homo \const{\graph}(\val, x, y))                             \\
                                                & \quad\Leftrightarrow\quad \tuple{x, y} \in \hat{\val}(\begin{tikzpicture}[baseline = -.5ex]
                                                                                                    \graph[grow right = 1.2cm, branch down = 6ex, nodes={mynode}]{
                                                                                                    {0/{}[draw, circle]}-!-{1/{}[draw, circle]}
                                                                                                    };
                                                                                                    \node[left = .5em of 0](l){};
                                                                                                    \node[right = .5em of 1](r){};
                                                                                                    \path (0) edge [draw = white, opacity = 0] node[pos= 0.5, inner sep = 1.5pt, font = \scriptsize, opacity = 1](a1){$\id$}(1);
                                                                                                    \graph[use existing nodes, edges={color=black, pos = .5, earrow}, edge quotes={fill=white, inner sep=1pt,font= \scriptsize}]{
                                                                                                           0 -- a1 -> 1;
                                                                                                           l -> 0; 1 -> r;
                                                                                                    };
                                                                                             \end{tikzpicture}) \tag{Def.\ of $\hat{\val}$} \\
                                                & \quad\Leftrightarrow\quad \tuple{x, y} \in \hat{\val}(\begin{tikzpicture}[baseline = -.5ex]
                                                                                                    \graph[grow right = 1.2cm, branch down = 6ex, nodes={mynode}]{
                                                                                                    {0/{}[draw, circle]}
                                                                                                    };
                                                                                                    \node[left = .5em of 0](l){};
                                                                                                    \node[right = .5em of 0](r){};
                                                                                                    \graph[use existing nodes, edges={color=black, pos = .5, earrow}, edge quotes={fill=white, inner sep=1pt,font= \scriptsize}]{
                                                                                                           l -> 0; 0 -> r;
                                                                                                    };
                                                                                             \end{tikzpicture}) \tag{$\id^{\const{\graph}(\val, x, y)}$ is the identity relation}                                 \\
                                                & \quad\Leftrightarrow\quad \tuple{x, y} \in \hat{\val}(\glang(\id)). \tag{Def.\ of $\glang$}
       \end{align*}

       Case $\term = \emp$:
       For every $\tuple{x, y} \in \hat{\val}(\top)$, we have
       \begin{align*}
              \tuple{x, y} \in \hat{\val}(\emp) & \Leftrightarrow \const{false}  \Leftrightarrow \tuple{x, y} \in \hat{\val}(\emptyset) \Leftrightarrow \tuple{x, y} \in \hat{\val}(\glang(\emp)). \tag{Def.\ of $\hat{\val}$ and $\glang$}
       \end{align*}

       Case $\term = \term[2] \compo \term[3]$:
       \begin{align*}
              \hat{\val}(\term[2] \compo \term[3]) & \quad=\quad \hat{\val}(\term[2]) \rcompo \hat{\val}(\term[3])  \tag{Def.\ of $\hat{\val}$}                                                                                                                     \\
                                                  & \quad=\quad \hat{\val}(\glang(\term[2])) \rcompo \hat{\val}(\glang(\term[3]))          \tag{IH}                                                                                                                \\
                                                  & \quad=\quad \bigcup_{\graph[1] \in \glang(\term[2])} \bigcup_{\graph[2] \in \glang(\term[3])} \hat{\val}(\graph[1]) \rcompo \hat{\val}(\graph[2])                 \tag{$\rcompo$ is distributive w.r.t.\ $\cup$} \\
                                                  & \quad=\quad \bigcup_{\graph[1] \in \glang(\term[2])} \bigcup_{\graph[2] \in \glang(\term[3])} \hat{\val}(\graph[1] \compo \graph[2]) \tag{\Cref{lemma: glang val compo}}                                      \\
                                                  & \quad=\quad \hat{\val}(\glang(\term[2] \compo \term[3])). \tag{Def.\ of $\glang$}
       \end{align*}

       Case $\term = \term[2] \cap \term[3]$:
       \begin{align*}
              \hat{\val}(\term[2] \cap \term[3]) & \quad=\quad \hat{\val}(\term[2]) \cap \hat{\val}(\term[3])                     \tag{Def.\ of $\hat{\val}$}                                                                                                  \\
                                                 & \quad=\quad \hat{\val}(\glang(\term[2])) \cap \hat{\val}(\glang(\term[3]))            \tag{IH}                                                                                                              \\
                                                 & \quad=\quad  \bigcup_{\graph[1] \in \glang(\term[2])} \bigcup_{\graph[2] \in \glang(\term[3])} (\hat{\val}(\graph[1]) \cap \hat{\val}(\graph[2]))                 \tag{$\cap$ is distributive w.r.t.\ $\cup$} \\
                                                 & \quad=\quad  \bigcup_{\graph[1] \in \glang(\term[2])} \bigcup_{\graph[2] \in \glang(\term[3])} \hat{\val}(\graph[1] \cap \graph[2]) \tag{\Cref{lemma: glang val cap}}                                       \\
                                                 & \quad=\quad \hat{\val}(\glang(\term[2] \cap \term[3])). \tag{Def.\ of $\glang$}
       \end{align*}

       Case $\term = \term[2] \union \term[3]$:
       \begin{align*}
              \hat{\val}(\term[2] \union \term[3]) & \quad=\quad \hat{\val}(\term[2]) \cup \hat{\val}(\term[3])                     \tag{Def.\ of $\hat{\val}$} \\
                                                 & \quad=\quad \hat{\val}(\glang(\term[2])) \cup \hat{\val}(\glang(\term[3]))            \tag{IH}             \\
                                                 & \quad=\quad \hat{\val}(\glang(\term[2]) \cup \glang(\term[3])) = \hat{\val}(\glang(\term[2] \union \term[3])). \tag{Def.\ of $\glang$}
       \end{align*}

       Case $\term = \term[2]^*$:
       \begin{align*}
              \hat{\val}(\term[2]^{*}) & \quad=\quad \bigcup_{n \ge 0} \hat{\val}(\underbrace{\term[2]\compo \dots \compo \term[2]}_{n \mbox{ times}}) = \bigcup_{n \ge 0} \underbrace{\hat{\val}(\term[2]) \compo \dots \compo \hat{\val}(\term[2])}_{n \mbox{ times}}  \tag{Def.\ of $\hat{\val}$}                                                                    \\
                                       & \quad=\quad \bigcup_{n \ge 0} \underbrace{\hat{\val}(\glang(\term[2])) \compo \dots \compo \hat{\val}(\glang(\term[2]))}_{n \mbox{ times}} \tag{IH}                                                                  \\
                                       & \quad=\quad \bigcup_{n \ge 0} \bigcup_{\graph_1, \dots, \graph_n \in  \glang(\term[2])} \hat{\val}(\graph_1) \rcompo \ldots \rcompo \hat{\val}(\graph_n)                         \tag{$\rcompo$ is distributive w.r.t.\ $\cup$} \\
                                       & \quad=\quad \bigcup_{n \ge 0} \bigcup_{\graph_1, \dots, \graph_n \in  \glang(\term[2])} \hat{\val}(\graph_1 \compo \ldots \compo \graph_n)  \tag{\Cref{lemma: glang val compo}}                                             \\
                                       & \quad=\quad \bigcup_{n \ge 0} \hat{\val}(\glang(\term[2]^n) = \hat{\val}(\glang(\term[2]^*)). \tag{Def.\ of $\glang$}
       \end{align*}
\end{proof}





% \section{Proof of \Cref{lemma: deduction theorem}}\label{section: lemma: deduction theorem}


% \begin{lemma}\label{lemma: restriction}
%        Let $\val \in \RSUB$ and $\tuple{x, y} \in \hat{\val}(\top)$.
%        Let $\val' \in \RSUB$ be the \kl{valuation} s.t.\ $\const{\graph}(\val', x, y)$ is the \kl{induced subgraph} of $\const{\graph}(\val, x, y)$ on the set $\set{z \mid \set{\tuple{x, z}, \tuple{z, y}} \subseteq \hat{\val}(\top)}$.
%        Then for all \kl{terms} $\term$, we have $\hat{\val}'(\term) = \hat{\val}(\term) \cap \hat{\val}'(\top)$.
% \end{lemma}
% \begin{proof}
%        By easy induction on $\term$.
% \end{proof}
% \begin{lemma}\label{lemma: top top}
%        Let $n \in \nat$ and let $\val$ be a \kl{valuation} of $\RSUB_{n}$.
%        Then we have $\hat{\val}(\term[3]) = \emptyset \Leftrightarrow \tuple{0, n} \not\in \hat{\val}(\top \term[3] \top)$
% \end{lemma}
% \begin{proof}
%        $\Rightarrow$:
%        Clear.
%        $\Leftarrow$:
%        Let $\tuple{i, j} \in \hat{\val}(\term[3])$.
%        By $\tuple{0, i}, \tuple{j, n} \in \hat{\val}(\top)$, we have $\tuple{0, n} \in \hat{\val}(\top \term[3] \top)$.
% \end{proof}

% \begin{proof}[Proof of \Cref{lemma: deduction theorem}]
%        We have
%        \begin{align*}
%                & \RSUB_{\Gamma \cup \set{\term[3] = \emp}} \not\models \term[1] \le \term[2]                                                                                                                                    \\
%                & \Leftrightarrow \exists n \in \nat, \val \in \RSUB_{n} \cap \RSUB_{\Gamma},\  \hat{\val}(\term[3]) = \emptyset \land \tuple{0, n} \in \hat{\val}(\term[1]) \setminus \hat{\val}(\term[2]) \tag{\Cref{lemma: restriction}} \\
%                & \Leftrightarrow \exists n \in \nat, \val \in \RSUB_{n} \cap \RSUB_{\Gamma},\  \tuple{0, n} \in \hat{\val}(\term[1]) \setminus \hat{\val}(\term[2] \cup \top \term[3] \top) \tag{\Cref{lemma: top top}}                    \\
%                & \Leftrightarrow \RSUB_{\Gamma} \not\models \term[1] \le \term[2] \cup \top \term[3] \top \tag*{(\Cref{lemma: restriction}) \qedhere}
%        \end{align*}
% \end{proof}


\fi